\documentclass[11pt,a4paper]{article}

% ============================================================
% PACKAGES
% ============================================================
\usepackage[utf8]{inputenc}
\usepackage[T1]{fontenc}
\usepackage[english]{babel}
\usepackage{amsmath,amssymb,amsthm,mathtools}
\usepackage{mathrsfs}
\usepackage{tikz-cd}
\usepackage{enumitem}
\usepackage{hyperref}
\usepackage{geometry}
\usepackage{microtype}

\geometry{margin=1in}

% ============================================================
% THEOREM ENVIRONMENTS
% ============================================================
\theoremstyle{definition}
\newtheorem{definition}{Definition}[section]
\newtheorem{axiom}[definition]{Axiom}
\newtheorem{example}[definition]{Example}
\newtheorem{remark}[definition]{Remark}
\newtheorem{convention}[definition]{Convention}

\theoremstyle{plain}
\newtheorem{proposition}[definition]{Proposition}
\newtheorem{theorem}[definition]{Theorem}
\newtheorem{lemma}[definition]{Lemma}
\newtheorem{corollary}[definition]{Corollary}

% ============================================================
% MACROS
% ============================================================
\newcommand{\Lam}{\Lambda}
\newcommand{\Sp}{\mathrm{Sp}}
\newcommand{\Ge}{\mathrm{Ge}}
\newcommand{\To}{\mathrm{To}}
\newcommand{\Al}{\mathrm{Al}}
\newcommand{\Nu}{\mathrm{Nu}}
\newcommand{\Fo}{\mathrm{Fo}}
\newcommand{\Cl}{\mathrm{Cl}}
\newcommand{\Corr}{\mathrm{Corr}}
\newcommand{\Inv}{\mathrm{Inv}}
\newcommand{\An}{\mathrm{An}}
\newcommand{\Ob}{\mathfrak{O}}
\newcommand{\Ecal}{\mathcal{E}}
\newcommand{\Ocal}{\mathcal{O}}
\newcommand{\Acal}{\mathcal{A}}
\newcommand{\philam}{\phi_{\lambda\mu}}
\newcommand{\rholam}{\rho_{\lambda}}
\newcommand{\Vmat}{\mathbf{V}}
\newcommand{\Jset}{\mathcal{J}}
\newcommand{\dd}{\,\mathrm{d}}

% ============================================================
% TITLE
% ============================================================
\title{%
  \textbf{A Multirepresentational Construction--Reconstruction\\
  Methodology for Mathematical Structures}\\[4pt]
  \large Part I: Foundations and Metamodel
}
\author{}
\date{}

\begin{document}
\maketitle

\begin{abstract}
\noindent
This document formalizes a methodology in which a single mathematical
object is represented through six interacting layers---spectral,
geometric, topological, algebraic, numerical, and formal---and
subsequently reconstructed through classical models, inverse
constructions, and analytic verification. Part~I establishes the
metamodel: the master object, the layer index set, the transition
morphisms, the compatibility conditions, the information invariants,
and the validation structure. The goal is to replace informal
narratives of ``stages'' by a compact, functorial, and checkable
formalism.
\end{abstract}

\tableofcontents

% ============================================================
\part{Foundations and Metamodel}
% ============================================================

% ------------------------------------------------------------
\section{Introduction}
% ------------------------------------------------------------

\subsection{Motivation}

Mathematical practice frequently moves between structurally different
descriptions of the same underlying object: an operator admits a
spectral decomposition, that decomposition generates a geometric
configuration, the configuration carries topological invariants,
those invariants admit algebraic codings, the algebraic codings admit
numerical approximations, and the whole edifice admits a formal
logical presentation. The prevailing habit is to treat these
descriptions as sequential steps.

This document adopts a different stance: the descriptions are treated
as \emph{projections of a single master object}, related by explicitly
declared morphisms, subject to commutativity and information-conservation
constraints.

\subsection{Problem Statement}

Given a mathematical object $\Ob$, we seek:
\begin{enumerate}[label=(\roman*)]
  \item a finite index set $\Lam$ of \emph{layers};
  \item for each $\lambda \in \Lam$, a projection
        $\pi_\lambda : \Ob \to \Ocal_\lambda$;
  \item for each ordered pair $\lambda \prec \mu$, a transition
        $\philam : \Ocal_\lambda \to \Ocal_\mu$;
  \item a compatibility criterion for the resulting diagram;
  \item a reconstruction functor from the formal layer back to
        classical models and corrected equations;
  \item a validation structure separating proof, computation, and
        conjecture.
\end{enumerate}

\subsection{Contributions}

The contributions of Part~I are:
\begin{itemize}
  \item a precise definition of the master object $\Ob$ and its
        layer decomposition (\S\ref{sec:notation});
  \item the definition of the transition morphisms and their
        compatibility axioms (\S\ref{sec:metamodel});
  \item a formal treatment of information loss and invariants
        (\S\ref{sec:compat});
  \item the introduction of the validation matrix $\Vmat$
        (\S\ref{sec:compat}).
\end{itemize}

\subsection{Organization}

Section~\ref{sec:notation} fixes notation. Section~\ref{sec:metamodel}
defines the six-layer metamodel. Section~\ref{sec:compat} treats
compatibility, invariants, and the validation matrix.

% ------------------------------------------------------------
\section{Notation and Conventions}\label{sec:notation}
% ------------------------------------------------------------

\subsection{The Master Object}

\begin{definition}[Master object]
Let $\Lam$ be a finite index set of \emph{layers}. A \emph{master
object} is a tuple
\[
  \Ob \;=\; \bigl(\Ocal_\lambda\bigr)_{\lambda \in \Lam},
\]
where each $\Ocal_\lambda$ is a mathematical structure (set, space,
algebra, category, or formal theory) associated with the layer
$\lambda$, together with a family of projection maps
\[
  \pi_\lambda : \Ob \longrightarrow \Ocal_\lambda,
  \qquad \lambda \in \Lam .
\]
\end{definition}

\begin{convention}[Standard layer set]
Unless otherwise stated,
\[
  \Lam \;=\; \{\Sp,\Ge,\To,\Al,\Nu,\Fo\},
\]
corresponding to the spectral, geometric, topological, algebraic,
numerical, and formal layers respectively.
\end{convention}

\begin{remark}
The master object $\Ob$ is not assumed to be a set-theoretic
primitive. It is a \emph{bookkeeping device} whose content is
exhausted by the data $(\Ocal_\lambda, \pi_\lambda)_{\lambda \in \Lam}$
and the transitions defined in \S\ref{sec:metamodel}.
\end{remark}

\subsection{Transition Morphisms}

\begin{definition}[Order relation on layers]
A binary relation $\prec$ on $\Lam$ is a \emph{construction order} if
it is irreflexive, transitive, and admits a linear extension
\[
  \Sp \prec \Ge \prec \To \prec \Al \prec \Nu \prec \Fo .
\]
\end{definition}

\begin{definition}[Transition morphisms]
For each pair $\lambda \prec \mu$ in $\Lam$, a \emph{transition
morphism} is a map
\[
  \philam : \Ocal_\lambda \longrightarrow \Ocal_\mu
\]
that carries structural data of layer $\lambda$ into structural data
of layer $\mu$.
\end{definition}

The resulting pipeline is
\begin{equation}\label{eq:pipeline}
  \Ocal_\Sp
  \xrightarrow{\phi_{\Sp\Ge}}
  \Ocal_\Ge
  \xrightarrow{\phi_{\Ge\To}}
  \Ocal_\To
  \xrightarrow{\phi_{\To\Al}}
  \Ocal_\Al
  \xrightarrow{\phi_{\Al\Nu}}
  \Ocal_\Nu
  \xrightarrow{\phi_{\Nu\Fo}}
  \Ocal_\Fo .
\end{equation}

\subsection{Reconstruction Operators}

\begin{definition}[Reconstruction operators]
The \emph{reconstruction operators} are maps
\[
  \rho_\Cl    : \Ocal_\Fo \to \Ecal_\Cl, \qquad
  \rho_\Inv   : \Ecal_\Cl \to \Ecal_\Corr, \qquad
  \rho_\An    : \Ecal_\Corr \to \Acal,
\]
where $\Ecal_\Cl$ denotes the classical-model layer, $\Ecal_\Corr$ the
corrected-model layer, and $\Acal$ the analytic-regularity layer.
\end{definition}

The reconstruction pipeline is
\begin{equation}\label{eq:reconstruction}
  \Ocal_\Fo
  \xrightarrow{\rho_\Cl}
  \Ecal_\Cl
  \xrightarrow{\rho_\Inv}
  \Ecal_\Corr
  \xrightarrow{\rho_\An}
  \Acal .
\end{equation}

\subsection{Discrete Families}

\begin{definition}[Discrete family]
A \emph{discrete family} of $\Ob$ is a family
\[
  \Ob^{(N)} \;=\; \bigl(\Ocal^{(N)}_\lambda\bigr)_{\lambda \in \Lam},
  \qquad N \in \mathbb{N},
\]
together with a family of \emph{resolution morphisms}
\[
  \iota_N : \Ob^{(N)} \to \Ob .
\]
\end{definition}

\begin{definition}[Discrete-to-continuum convergence]
Let $d_\lambda$ be a metric (or pseudometric) on $\Ocal_\lambda$, and
let $w_\lambda > 0$ be weights with $\sum_\lambda w_\lambda = 1$. The
\emph{layer-aggregated error} is
\begin{equation}\label{eq:error}
  \varepsilon_N
  \;=\;
  d_\Lam\bigl(\Ob^{(N)}, \Ob\bigr)
  \;=\;
  \sum_{\lambda \in \Lam} w_\lambda \,
  d_\lambda\bigl(\Ocal^{(N)}_\lambda, \Ocal_\lambda\bigr).
\end{equation}
We write $\Ob^{(N)} \to \Ob$ if $\varepsilon_N \to 0$ as
$N \to \infty$.
\end{definition}

\subsection{Validation Matrix}

\begin{definition}[Validation criteria]
Let
\[
  \Jset \;=\;
  \{\mathrm{Ex},\mathrm{Co},\mathrm{In},\mathrm{Nu},\mathrm{St},
    \mathrm{Cv},\mathrm{Re},\mathrm{Rev}\}
\]
denote the set of \emph{validation criteria}: existence, coherence,
invariance, numerical consistency, stability, convergence, regularity,
and reversibility, respectively.
\end{definition}

\begin{definition}[Validation matrix]
The \emph{validation matrix} is
\[
  \Vmat \;=\;
  \bigl(v_{\lambda j}\bigr)_{\lambda \in \Lam,\; j \in \Jset},
  \qquad
  v_{\lambda j} \in \{\mathrm{P},\mathrm{N},\mathrm{O},\mathrm{F}\},
\]
where the entries denote \emph{Proved}, \emph{Numerical},
\emph{Open}, and \emph{Failed}, respectively.
\end{definition}

\begin{remark}
The validation matrix distinguishes four epistemic statuses. In
particular, $\mathrm{P}$ (proved) and $\mathrm{N}$ (numerically
supported) must not be conflated: numerical evidence does not
constitute proof.
\end{remark}

% ------------------------------------------------------------
\section{The Metamodel}\label{sec:metamodel}
% ------------------------------------------------------------

\subsection{Six-Layer Representation}

\begin{definition}[Layer contents]
For a master object $\Ob = (\Ocal_\lambda)_{\lambda \in \Lam}$:
\begin{enumerate}[label=\textbf{(\roman*)}]
  \item \textbf{Spectral layer} $\Ocal_\Sp$: a spectral datum
        $(T, \sigma(T), E_T)$, where $T$ is a closed operator on a
        Hilbert space $H$, $\sigma(T) \subset \mathbb{C}$ is its
        spectrum, and $E_T$ is its spectral measure.
  \item \textbf{Geometric layer} $\Ocal_\Ge$: a geometric datum
        $(M, g, \nabla, R)$, where $M$ is a smooth manifold, $g$ a
        Riemannian metric, $\nabla$ the Levi-Civita connection, and
        $R$ the Riemann curvature tensor.
  \item \textbf{Topological layer} $\Ocal_\To$: a topological datum
        $(X, \{H_k\}_k, \{\pi_k\}_k, \chi)$, comprising a space $X$
        and its homology, homotopy, and Euler characteristic.
  \item \textbf{Algebraic layer} $\Ocal_\Al$: an algebraic datum
        $(A, G, R, M)$, comprising an algebra $A$, a group $G$, a
        ring $R$, and a module $M$ encoding the layer structure.
  \item \textbf{Numerical layer} $\Ocal_\Nu$: a numerical datum
        $(T_N, A_N, \{x_i\}_{i=1}^N, \varepsilon_N)$.
  \item \textbf{Formal layer} $\Ocal_\Fo$: a formal theory
        $(\mathcal{L}, \Gamma, \vdash)$, comprising a language
        $\mathcal{L}$, a set of axioms $\Gamma$, and a derivability
        relation $\vdash$.
\end{enumerate}
\end{definition}

\subsection{Compatibility Conditions}

\begin{axiom}[Functoriality of transitions]\label{ax:funct}
For every composable pair $\lambda \prec \mu \prec \nu$,
\[
  \phi_{\mu\nu} \circ \philam
  \;=\;
  \phi_{\lambda\nu} .
\]
\end{axiom}

\begin{axiom}[Projection coherence]\label{ax:proj}
For every $\lambda \in \Lam$ and every discrete family
$\Ob^{(N)} \to \Ob$,
\[
  \pi_\lambda \circ \iota_N
  \;=\;
  \pi^{(N)}_\lambda ,
\]
where $\pi^{(N)}_\lambda : \Ob^{(N)} \to \Ocal^{(N)}_\lambda$ is the
layer projection at resolution $N$.
\end{axiom}

\begin{axiom}[Reconstruction correspondence]\label{ax:corr}
There exists a \emph{correspondence limit}
\[
  \lim_{\epsilon \to 0} \rho_\Inv \circ \rho_\Cl
  \bigl(\Ocal_\Fo^{(\epsilon)}\bigr)
  \;=\;
  \Ecal_\Cl ,
\]
where $\Ocal_\Fo^{(\epsilon)}$ denotes the formal layer under an
$\epsilon$-perturbation of its axioms.
\end{axiom}

\subsection{Information Loss and Invariants}

\begin{definition}[Invariant of a layer]
An \emph{invariant} of layer $\lambda$ is a map
\[
  I_\lambda : \Ocal_\lambda \to \mathcal{V}_\lambda
\]
into a value object $\mathcal{V}_\lambda$, stable under the
\emph{admissible transformations} of $\Ocal_\lambda$:
\[
  U \sim U' \;\Longrightarrow\; I_\lambda(U) = I_\lambda(U').
\]
\end{definition}

\begin{definition}[Information loss]
The \emph{information loss} of the transition $\philam$ is the
subobject
\[
  \ker_\lambda(\philam)
  \;=\;
  \{ U \in \Ocal_\lambda : \philam(U) = \philam(U') \text{ for some } U' \neq U \}.
\]
The transition is \emph{information-preserving} if
$\ker_\lambda(\philam) = \varnothing$.
\end{definition}

\begin{proposition}[Invariant transport]\label{prop:inv-transport}
Let $I_\lambda$ be an invariant of $\Ocal_\lambda$. If $\philam$ is
information-preserving on the admissible class of $\Ocal_\lambda$,
then $I_\lambda$ descends to an invariant of $\Ocal_\mu$, i.e.\ there
exists $\tilde{I}_\mu$ such that
\[
  \tilde{I}_\mu \circ \philam = I_\lambda .
\]
\end{proposition}

\begin{proof}
Define $\tilde{I}_\mu$ on $\philam(\Ocal_\lambda)$ by
$\tilde{I}_\mu(\philam(U)) := I_\lambda(U)$; extend arbitrarily
elsewhere. Well-definedness follows from information preservation.
\end{proof}

\subsection{Validation Criteria}

\begin{definition}[Validation assignment]
A \emph{validation assignment} for layer $\lambda$ is a map
\[
  V_\lambda : \Jset \to \{\mathrm{P},\mathrm{N},\mathrm{O},\mathrm{F}\},
\]
such that $v_{\lambda j} = V_\lambda(j)$ recovers the row of $\Vmat$
associated with $\lambda$.
\end{definition}

\begin{definition}[Closed methodology]
The methodology is \emph{closed} if:
\begin{enumerate}[label=(\alph*)]
  \item every layer has at least one criterion marked $\mathrm{P}$;
  \item no criterion is marked $\mathrm{F}$ in any layer;
  \item the reversibility criterion $\mathrm{Rev}$ is marked
        $\mathrm{P}$ in the formal layer;
  \item the numerical criterion $\mathrm{Nu}$ is marked
        $\mathrm{N}$ or $\mathrm{P}$ in the numerical layer.
\end{enumerate}
\end{definition}

% ------------------------------------------------------------
\section{Compatibility and Commutativity}\label{sec:compat}
% ------------------------------------------------------------

\subsection{Commutative Diagrams}

The construction and reconstruction pipelines assemble into the
following diagram.

\begin{equation}\label{eq:bigdiagram}
\begin{tikzcd}[row sep=large, column sep=large]
  \Ocal_\Sp \arrow[r,"\phi_{\Sp\Ge}"] \arrow[d, "\pi^{(N)}_{\Sp}"']
    & \Ocal_\Ge \arrow[r,"\phi_{\Ge\To}"] \arrow[d, "\pi^{(N)}_{\Ge}"']
    & \Ocal_\To \arrow[r,"\phi_{\To\Al}"] \arrow[d, "\pi^{(N)}_{\To}"']
    & \Ocal_\Al \arrow[d, "\pi^{(N)}_{\Al}"] \\
  \Ocal^{(N)}_\Sp \arrow[r, "\phi^{(N)}_{\Sp\Ge}"']
    & \Ocal^{(N)}_\Ge \arrow[r, "\phi^{(N)}_{\Ge\To}"']
    & \Ocal^{(N)}_\To \arrow[r, "\phi^{(N)}_{\To\Al}"']
    & \Ocal^{(N)}_\Al
\end{tikzcd}
\end{equation}

\begin{definition}[Coherence of the diagram]
Diagram~\eqref{eq:bigdiagram} is \emph{coherent} if every square
commutes:
\[
  \pi^{(N)}_\mu \circ \philam
  \;=\;
  \phi^{(N)}_{\lambda\mu} \circ \pi^{(N)}_\lambda ,
  \qquad \lambda \prec \mu .
\]
\end{definition}

\begin{proposition}[Diagonal reduction]\label{prop:diag}
If Diagram~\eqref{eq:bigdiagram} is coherent and each
$\phi^{(N)}_{\lambda\mu} \to \philam$ as $N \to \infty$, then the
diagonal composite equals the horizontal composite in the limit:
\[
  \lim_{N \to \infty}
  \phi^{(N)}_{\Al\Nu} \circ \phi^{(N)}_{\To\Al} \circ \phi^{(N)}_{\Ge\To} \circ \phi^{(N)}_{\Sp\Ge}
  \;=\;
  \phi_{\Al\Nu} \circ \phi_{\To\Al} \circ \phi_{\Ge\To} \circ \phi_{\Sp\Ge} .
\]
\end{proposition}

\begin{proof}
Immediate from the coherence squares and continuity of composition
of morphisms in the target category.
\end{proof}

\subsection{Natural Transformations Between Layers}

\begin{definition}[Layer naturality]
Let $\mathcal{C}_\lambda$ and $\mathcal{C}_\mu$ be the categories
associated with layers $\lambda$ and $\mu$, and let
$F_\lambda : \mathcal{C}_\lambda \to \mathcal{C}$ and
$F_\mu : \mathcal{C}_\mu \to \mathcal{C}$ be realisation functors
into a common category $\mathcal{C}$. A \emph{layer natural
transformation} is a natural transformation
\[
  \eta : F_\lambda \Longrightarrow F_\mu \circ \philam .
\]
\end{definition}

\begin{proposition}[Invariant naturality]
If $I_\lambda$ is representable as a functor
$\mathcal{C}_\lambda \to \mathcal{V}_\lambda$ and $\philam$ admits a
right adjoint, then $I_\lambda$ extends naturally to layer $\mu$.
\end{proposition}

\subsection{Equivalence Versus Approximation}

\begin{definition}[Strict equivalence]
Layers $\lambda$ and $\mu$ are \emph{strictly equivalent}, written
$\Ocal_\lambda \simeq \Ocal_\mu$, if $\philam$ and $\phi_{\mu\lambda}$
are mutually inverse isomorphisms.
\end{definition}

\begin{definition}[Asymptotic equivalence]
Layers $\lambda$ and $\mu$ are \emph{asymptotically equivalent},
written $\Ocal_\lambda \sim_\infty \Ocal_\mu$, if there exists an
$\varepsilon_N \to 0$ and morphisms
\[
  \philam^{(N)} : \Ocal_\lambda \to \Ocal_\mu , \qquad
  \phi_{\mu\lambda}^{(N)} : \Ocal_\mu \to \Ocal_\lambda ,
\]
satisfying
\[
  d_\lambda\bigl(\phi_{\mu\lambda}^{(N)} \circ \philam^{(N)}, \mathrm{id}_{\Ocal_\lambda}\bigr)
  \;\leq\; \varepsilon_N ,
\]
and symmetrically.
\end{definition}

\begin{remark}
Strict equivalence is the exception; asymptotic equivalence is the
rule. The methodology therefore treats compatibility as a limiting
statement, not as an algebraic identity.
\end{remark}

\begin{theorem}[Closure of the methodology]\label{thm:closure}
If
\begin{enumerate}[label=(\roman*)]
  \item Diagram~\eqref{eq:bigdiagram} is coherent,
  \item $\Ob^{(N)} \to \Ob$ in the sense of~\eqref{eq:error},
  \item Axioms~\ref{ax:funct}, \ref{ax:proj}, \ref{ax:corr} hold,
\end{enumerate}
then the reconstruction pipeline~\eqref{eq:reconstruction} yields a
bidirectional characterization
\[
  \Ob \;\rightleftarrows\; \Ecal_\Corr .
\]
\end{theorem}

\begin{proof}[Sketch of proof]
Coherence provides the forward morphism from $\Ob$ to $\Ecal_\Fo$ via
the construction pipeline. Axiom~\ref{ax:corr} provides a
correspondence limit through $\rho_\Cl$ and $\rho_\Inv$, yielding
$\Ecal_\Corr$. Asymptotic equivalence between layers, via
Proposition~\ref{prop:diag}, supplies the reverse morphism. A
detailed proof requires exhibiting the reverse morphism explicitly for
the specific problem instance.
\end{proof}

\begin{remark}[Limits of the metamodel]
Theorem~\ref{thm:closure} is a \emph{structural} statement. It does
not, by itself, demonstrate that any particular mathematical object
admits a closed methodology. Verification must be carried out
layer by layer, and recorded in the validation matrix $\Vmat$.
\end{remark}

% ------------------------------------------------------------
\section*{Summary of Part I}
% ------------------------------------------------------------

Part~I has introduced:
\begin{itemize}
  \item the master object $\Ob$ and the standard layer set $\Lam$;
  \item the transition morphisms $\philam$ and their functoriality;
  \item the reconstruction operators $\rho_\Cl, \rho_\Inv, \rho_\An$;
  \item discrete families $\Ob^{(N)}$ and layer-aggregated error
        $\varepsilon_N$;
  \item the validation matrix $\Vmat$ and the four epistemic statuses
        $\{\mathrm{P},\mathrm{N},\mathrm{O},\mathrm{F}\}$;
  \item the compatibility axioms and the coherence of the big diagram;
  \item the distinction between strict and asymptotic equivalence;
  \item the closure theorem for the methodology.
\end{itemize}

Part~II will treat the construction pipeline layer by layer.

\documentclass[11pt,a4paper]{article}

% ============================================================
% PACKAGES
% ============================================================
\usepackage[utf8]{inputenc}
\usepackage[T1]{fontenc}
\usepackage[english]{babel}
\usepackage{amsmath,amssymb,amsthm,mathtools}
\usepackage{mathrsfs}
\usepackage{tikz-cd}
\usepackage{enumitem}
\usepackage{hyperref}
\usepackage{geometry}
\usepackage{microtype}

\geometry{margin=1in}

% ============================================================
% THEOREM ENVIRONMENTS
% ============================================================
\theoremstyle{definition}
\newtheorem{definition}{Definition}[section]
\newtheorem{axiom}[definition]{Axiom}
\newtheorem{example}[definition]{Example}
\newtheorem{remark}[definition]{Remark}
\newtheorem{convention}[definition]{Convention}
\newtheorem{construction}[definition]{Construction}

\theoremstyle{plain}
\newtheorem{proposition}[definition]{Proposition}
\newtheorem{theorem}[definition]{Theorem}
\newtheorem{lemma}[definition]{Lemma}
\newtheorem{corollary}[definition]{Corollary}

% ============================================================
% MACROS
% ============================================================
\newcommand{\Lam}{\Lambda}
\newcommand{\Sp}{\mathrm{Sp}}
\newcommand{\Ge}{\mathrm{Ge}}
\newcommand{\To}{\mathrm{To}}
\newcommand{\Al}{\mathrm{Al}}
\newcommand{\Nu}{\mathrm{Nu}}
\newcommand{\Fo}{\mathrm{Fo}}
\newcommand{\Cl}{\mathrm{Cl}}
\newcommand{\Corr}{\mathrm{Corr}}
\newcommand{\Inv}{\mathrm{Inv}}
\newcommand{\An}{\mathrm{An}}
\newcommand{\Ob}{\mathfrak{O}}
\newcommand{\Ecal}{\mathcal{E}}
\newcommand{\Ocal}{\mathcal{O}}
\newcommand{\Acal}{\mathcal{A}}
\newcommand{\Mcal}{\mathcal{M}}
\newcommand{\Hcal}{\mathcal{H}}
\newcommand{\philam}{\phi_{\lambda\mu}}
\newcommand{\rholam}{\rho_{\lambda}}
\newcommand{\Vmat}{\mathbf{V}}
\newcommand{\Jset}{\mathcal{J}}
\newcommand{\dd}{\,\mathrm{d}}
\newcommand{\op}[1]{\operatorname{#1}}
\newcommand{\spec}{\sigma}
\newcommand{\hilb}{\mathcal{H}}
\newcommand{\Cat}{\mathbf{Cat}}
\newcommand{\Set}{\mathbf{Set}}
\newcommand{\Vect}{\mathbf{Vect}}
\newcommand{\Grp}{\mathbf{Grp}}
\newcommand{\Rng}{\mathbf{Rng}}
\newcommand{\Top}{\mathbf{Top}}
\newcommand{\Man}{\mathbf{Man}}

% ============================================================
% TITLE
% ============================================================
\title{%
  \textbf{A Multirepresentational Construction--Reconstruction\\
  Methodology for Mathematical Structures}\\[4pt]
  \large Part II: The Construction Pipeline
}
\author{}
\date{}

\begin{document}
\maketitle

\begin{abstract}
\noindent
Part~II develops the construction pipeline of the multirepresentational
methodology. Starting from a master object $\Ob$, we construct, layer
by layer, the spectral datum, the geometric datum, the topological
datum, the algebraic datum, the numerical datum, and the formal datum.
Each layer is presented with: (i) its primitive objects; (ii) the
transition morphism that produces it from the previous layer;
(iii) the invariants it carries; and (iv) the validation criteria
that it must satisfy. The construction is treated as an ordered
composite of functors, and the deliverable of each layer is recorded
as an element of the corresponding projection $\Ocal_\lambda$.
\end{abstract}

\tableofcontents

% ============================================================
\part{The Construction Pipeline}
% ============================================================

\section{Overview of the Construction Pipeline}
\label{sec:overview}

Recall from Part~I the layer set
\[
  \Lam = \{\Sp, \Ge, \To, \Al, \Nu, \Fo\}
\]
and the construction order
\[
  \Sp \prec \Ge \prec \To \prec \Al \prec \Nu \prec \Fo .
\]
The construction pipeline is the composite
\begin{equation}\label{eq:pipeline-II}
  \Ocal_\Sp
  \xrightarrow{\phi_{\Sp\Ge}}
  \Ocal_\Ge
  \xrightarrow{\phi_{\Ge\To}}
  \Ocal_\To
  \xrightarrow{\phi_{\To\Al}}
  \Ocal_\Al
  \xrightarrow{\phi_{\Al\Nu}}
  \Ocal_\Nu
  \xrightarrow{\phi_{\Nu\Fo}}
  \Ocal_\Fo .
\end{equation}

\begin{convention}[Deliverable of a layer]
For each $\lambda \in \Lam$, the \emph{deliverable} of layer $\lambda$
is the pair
\[
  \Ocal_\lambda = (\mathcal{P}_\lambda, \mathcal{I}_\lambda),
\]
where $\mathcal{P}_\lambda$ is the set of primitive objects of the
layer and $\mathcal{I}_\lambda$ is the set of invariants carried by
the layer.
\end{convention}

\begin{remark}
Each section below follows the same structural template:
primitives, transition, invariants, and validation.
This uniformity is essential: it is what makes the layers comparable
and the validation matrix $\Vmat$ well-defined.
\end{remark}

% ------------------------------------------------------------
\section{Spectral Layer}
\label{sec:spectral}
% ------------------------------------------------------------

\subsection{Primitives}

\begin{definition}[Spectral datum]
A \emph{spectral datum} is a triple
\[
  \Ocal_\Sp = (T, \spec(T), E_T),
\]
where:
\begin{enumerate}[label=(\roman*)]
  \item $H$ is a separable complex Hilbert space;
  \item $T : D(T) \subset H \to H$ is a closed densely defined linear
        operator;
  \item $\spec(T) \subset \mathbb{C}$ is the spectrum of $T$;
  \item $E_T : \mathcal{B}(\spec(T)) \to \mathcal{B}(H)$ is the
        spectral measure associated with $T$, when $T$ is normal.
\end{enumerate}
\end{definition}

\begin{axiom}[Spectral realizability]
The operator $T$ is assumed normal, or self-adjoint, or essentially
self-adjoint on a dense domain $D(T)$.
\end{axiom}

Under this hypothesis, the spectral theorem yields
\begin{equation}\label{eq:spectral-theorem}
  T \;=\; \int_{\spec(T)} \lambda \dd E_T(\lambda),
\end{equation}
and the associated quadratic form
\[
  \langle T u, u \rangle
  \;=\;
  \int_{\spec(T)} \lambda \dd \mu_u(\lambda),
  \qquad
  \mu_u(\cdot) = \langle E_T(\cdot) u, u \rangle .
\]

\subsection{Transition}

\begin{definition}[Spectral--geometric transition]
The transition morphism
\[
  \phi_{\Sp\Ge} : \Ocal_\Sp \to \Ocal_\Ge
\]
is defined on eigenpairs by
\[
  \phi_{\Sp\Ge}(\lambda_i, v_i) \;=\; p_i \in M,
\]
where $M$ is a configuration space, manifold, or phase space on which
the spectral data are realised, and $p_i$ is the point or
configuration associated with the $i$-th mode.
\end{definition}

\begin{remark}
The transition $\phi_{\Sp\Ge}$ is not canonical: it depends on the
choice of realisation $M$. Different realisations produce different
geometric data and, consequently, different topological and algebraic
encodings. The choice of $M$ is itself part of the construction and
must be declared.
\end{remark}

\subsection{Invariants}

\begin{definition}[Spectral invariants]
The following are invariants of the spectral layer:
\begin{enumerate}[label=(\roman*)]
  \item the spectrum $\spec(T)$;
  \item the multiplicities
        $m(\lambda) = \dim \ker(T - \lambda I)$;
  \item the spectral measure $E_T$;
  \item the essential spectrum $\spec_{\mathrm{ess}}(T)$;
  \item the discrete spectrum $\spec_{\mathrm{disc}}(T)$;
  \item the spectral gap
        $\gamma(T) = \inf\{ |\lambda - \mu| : \lambda, \mu \in \spec(T), \lambda \neq \mu \}$.
\end{enumerate}
\end{definition}

\begin{proposition}[Unitary invariance]
If $U : H \to H'$ is unitary and $T' = U T U^{-1}$, then
\[
  \spec(T') = \spec(T), \qquad
  E_{T'}(B) = U E_T(B) U^{-1}.
\]
\end{proposition}

\subsection{Validation}

\begin{center}
\begin{tabular}{c|c|l}
  Criterion & Status & Comment \\
  \hline
  Existence & P & Spectral theorem \\
  Coherence & P & Functorial under unitary equivalence \\
  Invariance & P & Unitary invariance \\
  Numerical & N & Eigenvalue approximation \\
  Stability & P & Weyl's perturbation theorem \\
  Convergence & P & Spectral convergence \\
  Regularity & O & Depends on problem \\
  Reversibility & O & Inversion via spectral data \\
\end{tabular}
\end{center}

\begin{theorem}[Weyl's perturbation bound]
Let $T, T'$ be self-adjoint operators with $T' = T + V$ where
$V$ is bounded with $\|V\| = \epsilon$. Then for every $\lambda \in
\spec(T)$ there exists $\lambda' \in \spec(T')$ with
$|\lambda - \lambda'| \leq \epsilon$.
\end{theorem}

\begin{proof}
Standard; see Reed--Simon, \emph{Methods of Modern Mathematical
Physics}, Vol.~IV, Theorem~XIII.109.
\end{proof}

% ------------------------------------------------------------
\section{Geometric Layer}
\label{sec:geometric}
% ------------------------------------------------------------

\subsection{Primitives}

\begin{definition}[Geometric datum]
A \emph{geometric datum} is a tuple
\[
  \Ocal_\Ge = (M, g, \nabla, R),
\]
where:
\begin{enumerate}[label=(\roman*)]
  \item $M$ is a smooth finite-dimensional manifold;
  \item $g$ is a Riemannian (or pseudo-Riemannian) metric on $M$;
  \item $\nabla$ is the Levi-Civita connection of $g$;
  \item $R$ is the Riemann curvature tensor of $\nabla$.
\end{enumerate}
\end{definition}

\subsection{Transition}

\begin{definition}[Geometric--topological transition]
The transition morphism
\[
  \phi_{\Ge\To} : \Ocal_\Ge \to \Ocal_\To
\]
sends the geometric datum $(M, g, \nabla, R)$ to the underlying
topological datum $(X, \{H_k\}, \{\pi_k\}, \chi)$ with $X = M$
regarded as a topological space, forgetting the metric and connection.
\end{definition}

\subsection{Invariants}

\begin{definition}[Geometric invariants]
The geometric layer carries:
\begin{enumerate}[label=(\roman*)]
  \item the curvature tensor $R$ and its contractions
        $\op{Ric}$, $S$ (scalar curvature);
  \item the geodesic distance $d_g$;
  \item the volume form $\dd \mathrm{vol}_g$;
  \item the spectral data of the Laplace--Beltrami operator
        $\Delta_g$;
  \item the injectivity radius $\op{inj}(M, g)$;
  \item the sectional curvature $\op{sec}(g)$.
\end{enumerate}
\end{definition}

\begin{proposition}[Isometric invariance]
If $\varphi : (M, g) \to (M', g')$ is an isometry, then
\[
  R' = \varphi_* R, \qquad
  \op{Ric}' = \varphi_* \op{Ric}, \qquad
  S' = S \circ \varphi^{-1}.
\]
\end{proposition}

\subsection{Validation}

\begin{center}
\begin{tabular}{c|c|l}
  Criterion & Status & Comment \\
  \hline
  Existence & P & Smooth manifold theory \\
  Coherence & P & Natural under isometry \\
  Invariance & P & Isometric invariance \\
  Numerical & N & Discrete curvature, FEM \\
  Stability & O & Depends on curvature bounds \\
  Convergence & O & Cheeger--Gromov theory \\
  Regularity & P & Smoothness of $g$ \\
  Reversibility & O & Reconstructing $g$ from spectra \\
\end{tabular}
\end{center}

\begin{remark}
The transition from the spectral to the geometric layer is essentially
the \emph{inverse spectral problem}: reconstructing a metric from
spectral data of an associated operator (e.g.\ $\Delta_g$). This is a
classical ill-posed problem in the sense of Hadamard; see
Kac's problem on ``can one hear the shape of a drum?''.
\end{remark}

% ------------------------------------------------------------
\section{Topological Layer}
\label{sec:topological}
% ------------------------------------------------------------

\subsection{Primitives}

\begin{definition}[Topological datum]
A \emph{topological datum} is a tuple
\[
  \Ocal_\To = (X, \{H_k\}_k, \{\pi_k\}_k, \chi),
\]
where:
\begin{enumerate}[label=(\roman*)]
  \item $X$ is a topological space (typically a CW-complex or a
        manifold);
  \item $H_k = H_k(X; \mathbb{Z})$ is the $k$-th singular homology
        group;
  \item $\pi_k = \pi_k(X, x_0)$ is the $k$-th homotopy group;
  \item $\chi(X) = \sum_k (-1)^k \op{rk} H_k(X)$ is the Euler
        characteristic.
\end{enumerate}
\end{definition}

\subsection{Transition}

\begin{definition}[Topological--algebraic transition]
The transition morphism
\[
  \phi_{\To\Al} : \Ocal_\To \to \Ocal_\Al
\]
sends the topological datum to an algebraic datum by applying the
singular chain functor
\[
  C_\bullet : \Top \to \mathbf{Ch}(\mathbb{Z}),
  \qquad
  X \mapsto C_\bullet(X; \mathbb{Z}),
\]
and passing to homology:
\[
  H_\bullet(X) = H_\bullet(C_\bullet(X)).
\]
\end{definition}

\subsection{Invariants}

\begin{definition}[Topological invariants]
The topological layer carries:
\begin{enumerate}[label=(\roman*)]
  \item homology $H_k(X)$ and cohomology $H^k(X)$;
  \item homotopy groups $\pi_k(X)$;
  \item the Euler characteristic $\chi(X)$;
  \item characteristic classes (Stiefel--Whitney, Chern, Pontryagin);
  \item the fundamental group $\pi_1(X)$;
  \item homotopy type $[X]$.
\end{enumerate}
\end{definition}

\begin{theorem}[Homotopy invariance]
If $f : X \to Y$ is a homotopy equivalence, then
\[
  f_* : H_k(X) \xrightarrow{\sim} H_k(Y), \qquad
  f_* : \pi_k(X) \xrightarrow{\sim} \pi_k(Y), \qquad
  \chi(X) = \chi(Y).
\]
\end{theorem}

\subsection{Validation}

\begin{center}
\begin{tabular}{c|c|l}
  Criterion & Status & Comment \\
  \hline
  Existence & P & Singular (co)homology theory \\
  Coherence & P & Functoriality of $H_\bullet$ \\
  Invariance & P & Homotopy invariance \\
  Numerical & N & Persistent homology, algorithms \\
  Stability & P & Stability of persistence diagrams \\
  Convergence & P & Vietoris--Rips convergence \\
  Regularity & P & CW-complex assumptions \\
  Reversibility & O & Reconstructing $X$ from $H_\bullet$ \\
\end{tabular}
\end{center}

\begin{remark}
The functor $H_k$ is a topological invariant of $X$ but not a
complete invariant: non-homotopy-equivalent spaces may share
homology. The passage from $\Ocal_\Ge$ to $\Ocal_\To$ therefore
entails a controlled information loss, in the sense of
Definition~1.9 of Part~I.
\end{remark}

% ------------------------------------------------------------
\section{Algebraic Layer}
\label{sec:algebraic}
% ------------------------------------------------------------

\subsection{Primitives}

\begin{definition}[Algebraic datum]
An \emph{algebraic datum} is a tuple
\[
  \Ocal_\Al = (A, G, R, \Mcal),
\]
where:
\begin{enumerate}[label=(\roman*)]
  \item $A$ is an associative (or Lie) algebra over a base ring;
  \item $G$ is a group encoding the symmetries of the previous
        layers;
  \item $R$ is a ring carrying the arithmetic structure of the
        invariants;
  \item $\Mcal$ is a module over $R$ on which $A$ acts.
\end{enumerate}
\end{definition}

\subsection{Transition}

\begin{definition}[Algebraic--numerical transition]
The transition morphism
\[
  \phi_{\Al\Nu} : \Ocal_\Al \to \Ocal_\Nu
\]
chooses a finite-dimensional truncation
\[
  A \;\longrightarrow\; A_N \in \mathrm{Mat}_{N \times N}(R),
\]
compatible with a filtration
\[
  A_1 \subset A_2 \subset \cdots \subset A_N \subset \cdots \subset A .
\]
\end{definition}

\subsection{Invariants}

\begin{definition}[Algebraic invariants]
The algebraic layer carries:
\begin{enumerate}[label=(\roman*)]
  \item the isomorphism class of $A$;
  \item the group $G$ and its representation theory;
  \item the ring structure of $R$;
  \item the module structure of $\Mcal$;
  \item the Hochschild homology/cohomology of $A$;
  \item the $K$-theory of $R$.
\end{enumerate}
\end{definition}

\begin{proposition}[Functoriality]
Every layer transition $\philam$ factors through the algebraic layer
whenever the target category admits a forgetful functor to
$\mathbf{Alg}$.
\end{proposition}

\subsection{Validation}

\begin{center}
\begin{tabular}{c|c|l}
  Criterion & Status & Comment \\
  \hline
  Existence & P & Standard algebraic constructions \\
  Coherence & P & Functoriality \\
  Invariance & P & Isomorphism invariance \\
  Numerical & N & Sparse linear algebra \\
  Stability & O & Conditioning of truncations \\
  Convergence & O & Depends on filtration \\
  Regularity & P & Algebraic identities exact \\
  Reversibility & O & Inverting the functor \\
\end{tabular}
\end{center}

% ------------------------------------------------------------
\section{Numerical Layer}
\label{sec:numerical}
% ------------------------------------------------------------

\subsection{Primitives}

\begin{definition}[Numerical datum]
A \emph{numerical datum} is a tuple
\[
  \Ocal_\Nu = (T_N, A_N, \{x_i\}_{i=1}^N, \varepsilon_N),
\]
where:
\begin{enumerate}[label=(\roman*)]
  \item $N \in \mathbb{N}$ is the resolution parameter;
  \item $T_N \in \mathrm{Mat}_{N \times N}(\mathbb{C})$ is a
        finite-dimensional approximation of $T$;
  \item $A_N$ is a finite-dimensional truncation of $A$;
  \item $\{x_i\}_{i=1}^N$ is a discretisation grid;
  \item $\varepsilon_N$ is the discretisation error.
\end{enumerate}
\end{definition}

\subsection{Transition}

\begin{definition}[Numerical--formal transition]
The transition morphism
\[
  \phi_{\Nu\Fo} : \Ocal_\Nu \to \Ocal_\Fo
\]
encodes the numerical datum into a formal language by declaring the
discretisation scheme as a set of axioms $\Gamma_N$ and the
approximation as a term $t_N$ provable in $\Gamma_N$:
\[
  \Gamma_N \vdash t_N : T_N .
\]
\end{definition}

\subsection{Discretisation Scheme}

\begin{definition}[Consistency]
The discretisation $\{T_N\}_{N \geq 1}$ is \emph{consistent} with $T$
if
\[
  \|T_N \pi_N - \pi_N T\| \;\longrightarrow\; 0,
  \qquad N \to \infty,
\]
where $\pi_N : H \to H_N$ is the orthogonal projection onto the
finite-dimensional subspace $H_N$.
\end{definition}

\begin{definition}[Stability]
The scheme is \emph{stable} if there exists $C > 0$ such that
\[
  \|T_N^{-1}\| \;\leq\; C
  \qquad \forall N \in \mathbb{N}.
\]
\end{definition}

\begin{theorem}[Lax--Richtmyer equivalence]\label{thm:lax-rich}
For a consistent discretisation of a well-posed linear problem,
\[
  \text{stability} \;\Longleftrightarrow\; \text{convergence}.
\]
\end{theorem}

\subsection{Invariants}

\begin{definition}[Numerical invariants]
The numerical layer carries:
\begin{enumerate}[label=(\roman*)]
  \item the discretisation error $\varepsilon_N$;
  \item the condition number $\kappa(T_N) = \|T_N\| \|T_N^{-1}\|$;
  \item the convergence rate $p$ such that
        $\varepsilon_N = \mathcal{O}(N^{-p})$;
  \item the stability constant $C$.
\end{enumerate}
\end{definition}

\subsection{Validation}

\begin{center}
\begin{tabular}{c|c|l}
  Criterion & Status & Comment \\
  \hline
  Existence & P & Finite-dimensional approximation \\
  Coherence & P & Projection consistency \\
  Invariance & P & Spectral invariance up to $\varepsilon_N$ \\
  Numerical & N & Simulation \\
  Stability & P & Lax--Richtmyer \\
  Convergence & P & Lax--Richtmyer \\
  Regularity & O & Depends on scheme \\
  Reversibility & O & Numerical inversion \\
\end{tabular}
\end{center}

\begin{remark}
The numerical layer is the only layer in which results can be
generated by a machine. For this reason it plays a privileged role in
\emph{structure discovery}---but not in \emph{structure proof}. The
distinction is captured by the validation matrix $\Vmat$, where
numerical entries are marked $\mathrm{N}$ rather than
$\mathrm{P}$.
\end{remark}

% ------------------------------------------------------------
\section{Formal Layer}
\label{sec:formal}
% ------------------------------------------------------------

\subsection{Primitives}

\begin{definition}[Formal datum]
A \emph{formal datum} is a triple
\[
  \Ocal_\Fo = (\mathcal{L}, \Gamma, \vdash),
\]
where:
\begin{enumerate}[label=(\roman*)]
  \item $\mathcal{L}$ is a first-order (or higher-order) language;
  \item $\Gamma \subseteq \mathrm{Sent}(\mathcal{L})$ is a set of
        axioms;
  \item $\vdash$ is the derivability relation of a sound and complete
        proof system (e.g.\ natural deduction, sequent calculus).
\end{enumerate}
\end{definition}

\subsection{Transition}

\begin{definition}[Formal--classical transition]
The transition morphism
\[
  \rho_\Cl : \Ocal_\Fo \to \Ecal_\Cl
\]
extracts from the formal theory $\Gamma$ a classical model
$\Ecal_\Cl$, defined as the class of structures $\Mcal$ satisfying
\[
  \Mcal \models \Gamma .
\]
\end{definition}

\subsection{Derivation Structure}

\begin{definition}[Derivation tree]
A \emph{derivation tree} in $\Gamma$ is a finite tree whose leaves
are axioms in $\Gamma$ and whose internal nodes are instances of
inference rules, terminating in a conclusion
\[
  \Gamma \vdash \varphi .
\]
\end{definition}

\begin{definition}[Layer formalisation]
The formal layer is the terminal layer of the construction pipeline
if and only if the following holds for every layer
$\lambda \in \Lam$:
\[
  \Gamma \vdash \varphi_\lambda,
\]
where $\varphi_\lambda$ is the sentence of $\mathcal{L}$ asserting the
existence and properties of the invariant $I_\lambda$.
\end{definition}

\subsection{Invariants}

\begin{definition}[Formal invariants]
The formal layer carries:
\begin{enumerate}[label=(\roman*)]
  \item the derivability relation $\vdash$;
  \item the consistency of $\Gamma$, i.e.\
        $\Gamma \nvdash \bot$;
  \item the completeness of the theory with respect to its intended
        model;
  \item the decidability (or undecidability) of the theory;
  \item the proof-theoretic ordinal of $\Gamma$.
\end{enumerate}
\end{definition}

\subsection{Validation}

\begin{center}
\begin{tabular}{c|c|l}
  Criterion & Status & Comment \\
  \hline
  Existence & P & By construction \\
  Coherence & P & Soundness of $\vdash$ \\
  Invariance & P & Logical invariance \\
  Numerical & O & Not applicable directly \\
  Stability & P & Under extension of $\Gamma$ \\
  Convergence & P & Completeness theorem \\
  Regularity & P & Proof-theoretic regularity \\
  Reversibility & P & Interpreting classical model \\
\end{tabular}
\end{center}

\begin{theorem}[Completeness of the formal layer]\label{thm:completeness}
Let $\Gamma$ be a consistent set of sentences of $\mathcal{L}$. Then
\[
  \Gamma \vdash \varphi
  \;\Longleftrightarrow\;
  \Gamma \models \varphi .
\]
\end{theorem}

\begin{proof}
G\"odel's completeness theorem (see Ebbinghaus--Flum--Thomas,
\emph{Mathematical Logic}, GTM~291, Theorem~X.2.1).
\end{proof}

\begin{remark}
Completeness of the formal layer is a syntactic-semantic equivalence
within $\mathcal{L}$. It does not establish that the axiomatisation
$\Gamma$ captures the intended mathematical object $\Ob$; this is a
question of \emph{soundness of formalisation} and lies outside the
scope of the completeness theorem.
\end{remark}

% ------------------------------------------------------------
\section{Composition of the Pipeline}
\label{sec:composition}
% ------------------------------------------------------------

\subsection{The Composite Functor}

\begin{definition}[Construction functor]
The \emph{construction functor} is the composite
\[
  \mathbf{C}
  \;=\;
  \phi_{\Nu\Fo}
  \circ
  \phi_{\Al\Nu}
  \circ
  \phi_{\To\Al}
  \circ
  \phi_{\Ge\To}
  \circ
  \phi_{\Sp\Ge}
  :
  \Ocal_\Sp \to \Ocal_\Fo .
\]
\end{definition}

\begin{proposition}[Functoriality]
The construction functor $\mathbf{C}$ is functorial in the sense of
Axiom~1.3 of Part~I:
\[
  \mathbf{C}(\philam \circ \phi_{\kappa\lambda})
  =
  \mathbf{C}(\philam) \circ \mathbf{C}(\phi_{\kappa\lambda})
\]
for every composable chain $\kappa \prec \lambda \prec \mu$ in $\Lam$.
\end{proposition}

\subsection{Information Flow}

\begin{definition}[Information flow diagram]
The information flow through the pipeline is the diagram
\[
\begin{tikzcd}[row sep=large, column sep=large]
  \Ocal_\Sp \arrow[r, "\phi_{\Sp\Ge}"] \arrow[rrrr, bend left=20, "\mathbf{C}_{\Sp\Fo}"]
    & \Ocal_\Ge \arrow[r, "\phi_{\Ge\To}"]
    & \Ocal_\To \arrow[r, "\phi_{\To\Al}"]
    & \Ocal_\Al \arrow[r, "\phi_{\Al\Nu}"]
    & \Ocal_\Nu \arrow[r, "\phi_{\Nu\Fo}"]
    & \Ocal_\Fo
\end{tikzcd}
\]
\end{definition}

\begin{proposition}[Cumulative information loss]
The total information loss along the pipeline is the composite
kernel
\[
  \ker(\mathbf{C})
  =
  \bigcup_{\lambda \prec \Fo} \ker_\lambda(\philam)
  \;\cup\;
  \ker_{\Sp}(\mathbf{C}),
\]
in the sense of Definition~1.9 of Part~I.
\end{proposition}

\subsection{Deliverables}

\begin{definition}[Deliverable chain]
The \emph{deliverable chain} of the construction pipeline is the
sequence
\[
  \Ocal_\Sp
  \;\longrightarrow\;
  \Ocal_\Ge
  \;\longrightarrow\;
  \Ocal_\To
  \;\longrightarrow\;
  \Ocal_\Al
  \;\longrightarrow\;
  \Ocal_\Nu
  \;\longrightarrow\;
  \Ocal_\Fo ,
\]
in which each arrow is a transition morphism in the sense of
Part~I, Definition~1.4.
\end{definition}

% ------------------------------------------------------------
\section*{Summary of Part II}
% ------------------------------------------------------------

Part~II has constructed, layer by layer, the six data of the
construction pipeline:
\[
  \Ocal_\Sp,\ \Ocal_\Ge,\ \Ocal_\To,\ \Ocal_\Al,\ \Ocal_\Nu,\ \Ocal_\Fo.
\]
Each layer was presented with:
\begin{itemize}
  \item its primitive objects (Definitions~2.2, 3.2, 4.2, 5.2, 6.2,
        7.2);
  \item its transition morphism from the previous layer
        (Definitions~2.4, 3.4, 4.4, 5.4, 6.4, 7.4);
  \item its invariants (Definitions~2.5, 3.5, 4.5, 5.5, 6.5, 7.5);
  \item its validation criteria.
\end{itemize}
The composite of the pipeline defines the construction functor
$\mathbf{C} : \Ocal_\Sp \to \Ocal_\Fo$, functorial by
Axiom~1.3 of Part~I.

Part~III will treat the reverse direction: reconstruction, inverse
problems, correction operators, and the bidirectional
characterisation.

\documentclass[11pt,a4paper]{article}

% ============================================================
% PACKAGES
% ============================================================
\usepackage[utf8]{inputenc}
\usepackage[T1]{fontenc}
\usepackage[english]{babel}
\usepackage{amsmath,amssymb,amsthm,mathtools}
\usepackage{mathrsfs}
\usepackage{tikz-cd}
\usepackage{enumitem}
\usepackage{hyperref}
\usepackage{geometry}
\usepackage{microtype}

\geometry{margin=1in}

% ============================================================
% THEOREM ENVIRONMENTS
% ============================================================
\theoremstyle{definition}
\newtheorem{definition}{Definition}[section]
\newtheorem{axiom}[definition]{Axiom}
\newtheorem{example}[definition]{Example}
\newtheorem{remark}[definition]{Remark}
\newtheorem{convention}[definition]{Convention}
\newtheorem{construction}[definition]{Construction}

\theoremstyle{plain}
\newtheorem{proposition}[definition]{Proposition}
\newtheorem{theorem}[definition]{Theorem}
\newtheorem{lemma}[definition]{Lemma}
\newtheorem{corollary}[definition]{Corollary}

% ============================================================
% MACROS
% ============================================================
\newcommand{\Lam}{\Lambda}
\newcommand{\Sp}{\mathrm{Sp}}
\newcommand{\Ge}{\mathrm{Ge}}
\newcommand{\To}{\mathrm{To}}
\newcommand{\Al}{\mathrm{Al}}
\newcommand{\Nu}{\mathrm{Nu}}
\newcommand{\Fo}{\mathrm{Fo}}
\newcommand{\Cl}{\mathrm{Cl}}
\newcommand{\Corr}{\mathrm{Corr}}
\newcommand{\Inv}{\mathrm{Inv}}
\newcommand{\An}{\mathrm{An}}
\newcommand{\Ob}{\mathfrak{O}}
\newcommand{\Ecal}{\mathcal{E}}
\newcommand{\Ocal}{\mathcal{O}}
\newcommand{\Acal}{\mathcal{A}}
\newcommand{\Mcal}{\mathcal{M}}
\newcommand{\Hcal}{\mathcal{H}}
\newcommand{\Dcal}{\mathcal{D}}
\newcommand{\philam}{\phi_{\lambda\mu}}
\newcommand{\rholam}{\rho_{\lambda}}
\newcommand{\Vmat}{\mathbf{V}}
\newcommand{\Jset}{\mathcal{J}}
\newcommand{\dd}{\,\mathrm{d}}
\newcommand{\op}[1]{\operatorname{#1}}
\newcommand{\spec}{\sigma}
\newcommand{\hilb}{\mathcal{H}}
\newcommand{\Cat}{\mathbf{Cat}}
\newcommand{\Set}{\mathbf{Set}}
\newcommand{\Vect}{\mathbf{Vect}}
\newcommand{\Grp}{\mathbf{Grp}}
\newcommand{\Rng}{\mathbf{Rng}}
\newcommand{\Top}{\mathbf{Top}}
\newcommand{\Man}{\mathbf{Man}}
\newcommand{\Norm}[1]{\left\lVert #1 \right\rVert}
\newcommand{\Abs}[1]{\left\lvert #1 \right\rvert}
\newcommand{\ip}[2]{\left\langle #1, #2 \right\rangle}

% ============================================================
% TITLE
% ============================================================
\title{%
  \textbf{A Multirepresentational Construction--Reconstruction\\
  Methodology for Mathematical Structures}\\[4pt]
  \large Part III: Reconstruction, Inversion, and Correction
}
\author{}
\date{}

\begin{document}
\maketitle

\begin{abstract}
\noindent
Part~III develops the reverse direction of the multirepresentational
methodology. Given the formal layer $\Ocal_\Fo$ produced by the
construction pipeline of Part~II, we reconstruct the classical model
$\Ecal_\Cl$, formulate the inverse problem of recovering the
generating structure from its formal encoding, and derive a
correction operator $\Delta L$ whose effect is measured by a
correspondence limit. The central objects are the reconstruction
operators $\rho_\Cl, \rho_\Inv, \rho_\An$ introduced in Part~I, the
classical-model layer $\Ecal_\Cl$, the corrected-model layer
$\Ecal_\Corr$, and the bidirectional characterisation
$\Ob \rightleftarrows \Ecal_\Corr$. We also treat well-posedness,
regularisation, and the stability conditions required for a
meaningful inversion.
\end{abstract}

\tableofcontents

% ============================================================
\part{Reconstruction, Inversion, and Correction}
% ============================================================

\section{Overview of the Reconstruction Pipeline}
\label{sec:overview-III}

Recall from Part~I the reconstruction operators
\[
  \rho_\Cl    : \Ocal_\Fo \to \Ecal_\Cl, \qquad
  \rho_\Inv   : \Ecal_\Cl \to \Ecal_\Corr, \qquad
  \rho_\An    : \Ecal_\Corr \to \Acal,
\]
and the reconstruction pipeline
\begin{equation}\label{eq:reconstruction-III}
  \Ocal_\Fo
  \xrightarrow{\rho_\Cl}
  \Ecal_\Cl
  \xrightarrow{\rho_\Inv}
  \Ecal_\Corr
  \xrightarrow{\rho_\An}
  \Acal .
\end{equation}

Part~III is organised as follows.
Section~\ref{sec:classical} treats the extraction of classical
models. Section~\ref{sec:inverse} formulates the inverse problem.
Section~\ref{sec:correction} defines the correction operator and the
correspondence limit. Section~\ref{sec:bidirectional} establishes the
bidirectional characterisation $\Ob \rightleftarrows \Ecal_\Corr$.
Section~\ref{sec:wellposedness} treats well-posedness, stability, and
regularisation.

\begin{remark}
The reconstruction pipeline is \emph{not} the formal inverse of the
construction pipeline. The latter is a sequence of forgetful and
realising functors; the former is a sequence of solution and
selection operators. This asymmetry is the origin of the
well-posedness questions treated in \S\ref{sec:wellposedness}.
\end{remark}

% ------------------------------------------------------------
\section{Classical Models}
\label{sec:classical}
% ------------------------------------------------------------

\subsection{The Classical-Model Layer}

\begin{definition}[Classical-model layer]\label{def:classical-layer}
The \emph{classical-model layer} is the class
\[
  \Ecal_\Cl
  \;=\;
  \bigl\{ \Mcal : \Mcal \models \Gamma \bigr\},
\]
where $\Gamma$ is the axiom set of the formal layer $\Ocal_\Fo$ and
$\models$ is the satisfaction relation of the underlying logic
$\mathcal{L}$.
\end{definition}

\begin{definition}[Distinguished classical model]
A \emph{distinguished classical model} is a pair
\[
  \Mcal_\star = (\Omega_\star, \mathcal{E}_\star)
\]
where $\Omega_\star$ is a domain (open subset of $\mathbb{R}^n$, a
manifold, a function space, or a probability space) and
$\mathcal{E}_\star$ is a classical equation
\[
  \mathcal{E}_\star(u) = 0, \qquad u \in \Omega_\star .
\]
\end{definition}

\subsection{Extraction}

\begin{construction}[Classical extraction]\label{constr:classical}
Given $\Ocal_\Fo = (\mathcal{L}, \Gamma, \vdash)$, the extraction
operator $\rho_\Cl$ proceeds as follows:
\begin{enumerate}[label=(\roman*)]
  \item select the distinguished symbols of $\mathcal{L}$ that
        encode operators, domains, and boundary conditions;
  \item interpret each selected symbol in $\Omega_\star$;
  \item collect the resulting equations into a system
        $\mathcal{E}_\star(u) = 0$;
  \item verify $\Mcal_\star \models \Gamma$.
\end{enumerate}
The output is the distinguished classical model $\Mcal_\star$.
\end{construction}

\begin{remark}
The extraction is not unique: distinct choices of distinguished
symbols may yield distinct classical models. When this occurs, the
selection is part of the reconstruction and must be declared. This
parallels the non-canonicity of $\phi_{\Sp\Ge}$ in
Part~II, \S\ref{sec:spectral}.
\end{remark}

\subsection{Identification of Known Theories}

\begin{definition}[Classical containment]
A classical equation $\mathcal{E}_0(u) = 0$ is \emph{contained} in
$\Ecal_\Cl$, written $\mathcal{E}_0 \trianglelefteq \Ecal_\Cl$, if
there exists $\Mcal \in \Ecal_\Cl$ such that
\[
  \mathcal{E}_\Mcal = \mathcal{E}_0 .
\]
\end{definition}

\begin{proposition}[Recognition criterion]\label{prop:recognition}
If $\mathcal{E}_0 \trianglelefteq \Ecal_\Cl$ then the extracted
model coincides with a known theory and no correction is required at
the level of $\mathcal{E}_0$.
\end{proposition}

\begin{proof}
By definition of containment, $\mathcal{E}_0$ satisfies the axioms
$\Gamma$ and hence is a fixed point of the extraction. The correction
operator (Definition~\ref{def:correction}) therefore acts trivially
on $\mathcal{E}_0$.
\end{proof}

\subsection{Correspondence Limits}

\begin{definition}[Correspondence limit]
Let $\{ L_\epsilon \}_{\epsilon > 0}$ be a family of operators
parametrised by $\epsilon$. A \emph{correspondence limit} is a limit
\[
  L_\star
  \;=\;
  \lim_{\epsilon \to 0^+} L_\epsilon ,
\]
in a topology $\tau$ specified on the operator space, together with
an identification
\[
  L_\star = L_\Cl ,
\]
where $L_\Cl$ is the operator of the distinguished classical model.
\end{definition}

\begin{remark}
The topology $\tau$ is not a technicality: strong, weak, norm,
and resolvent topologies yield different correspondence statements.
The choice of $\tau$ is part of the reconstruction and enters the
validation matrix $\Vmat$ at the entry $(\Fo, \mathrm{Cv})$.
\end{remark}

% ------------------------------------------------------------
\section{The Inverse Problem}
\label{sec:inverse}
% ------------------------------------------------------------

\subsection{Structural Inversion}

\begin{definition}[Structural inverse problem]\label{def:inverse-problem}
Given a classical equation
\[
  \mathcal{E}(u) = 0, \qquad u \in \Omega,
\]
and a target structure $\mathcal{P}$ on the domain $\Omega$, the
\emph{structural inverse problem} consists in finding an operator
$L^\ast$ such that
\[
  \mathcal{E}(u) = 0
  \;\Longleftrightarrow\;
  L^\ast u = \lambda u \quad \text{(or } L^\ast u = f \text{)}
\]
and such that $L^\ast$ realises $\mathcal{P}$.
\end{definition}

\begin{definition}[Inversion operator]
The \emph{inversion operator} is the map
\[
  \rho_\Inv : \Ecal_\Cl \to \Ecal_\Corr
\]
that assigns to a classical model $\Mcal_\star$ an operator
$L^\ast$ solving the structural inverse problem for $\Mcal_\star$.
\end{definition}

\subsection{The Inverse Spectral Paradigm}

\begin{example}[Inverse spectral problem]
Let $\Omega \subset \mathbb{R}^n$ be a bounded domain with smooth
boundary, and let
\[
  \mathcal{E}(u) = -\Delta u - \lambda u = 0, \qquad u|_{\partial \Omega} = 0 .
\]
The structural inverse problem is:
\emph{given the spectrum $\{\lambda_k\}_{k \geq 1}$ of $-\Delta$,
recover the domain $\Omega$ (up to isometry)}.
This is the classical problem of Kac, \emph{Can one hear the shape of
a drum?}.
\end{example}

\begin{remark}
The inverse spectral problem is \emph{not} uniquely solvable in
general: there exist isospectral but non-isometric domains
(Gordon--Webb--Wolpert, 1992). The failure of injectivity of
$\rho_\Inv$ is thus intrinsic, not an artefact of the method.
\end{remark}

\subsection{The Inverse Construction Problem}

\begin{definition}[Structure-to-equation problem]\label{def:structure-to-eq}
Given the deliverable chain of Part~II,
\[
  \Ocal_\Sp, \Ocal_\Ge, \Ocal_\To, \Ocal_\Al,
\]
a \emph{structure-to-equation problem} consists in finding an
operator $L^\ast$ such that the following compatibility holds:
\[
  \spec(L^\ast) \cong \Ocal_\Sp, \qquad
  (M, g) = \Ocal_\Ge, \qquad
  H_\bullet(M) = \Ocal_\To, \qquad
  A(M) = \Ocal_\Al .
\]
\end{definition}

\begin{proposition}[Non-uniqueness of structure-to-equation]\label{prop:nonunique}
In general, the structure-to-equation problem admits more than one
solution. The set of solutions is
\[
  \mathcal{S} = \rho_\Inv^{-1}(\Ecal_\Corr) ,
\]
and $\rho_\Inv$ is injective if and only if $\mathcal{S}$ is a
singleton.
\end{proposition}

\begin{proof}
Immediate from the definition of the preimage and injectivity.
\end{proof}

\subsection{The Correction Operator}

\begin{definition}[Delta operator]\label{def:correction}
Let $L_\Cl$ be the operator of the distinguished classical model and
let $L^\ast = \rho_\Inv(\Mcal_\star)$. The \emph{delta operator} is
\[
  \Delta L
  \;=\;
  L^\ast - L_\Cl .
\]
The \emph{corrected operator} is
\[
  L_\Corr
  \;=\;
  L_\Cl + \Delta L .
\]
\end{definition}

\begin{remark}
The delta operator $\Delta L$ is not assumed to be small. Its
magnitude is a diagnostic quantity, not a perturbative hypothesis.
The interpretation of $\Delta L$ is deferred to
\S\ref{sec:wellposedness}.
\end{remark}

\subsection{Layerwise Decomposition of the Correction}

\begin{definition}[Layerwise decomposition]
The delta operator admits a decomposition
\[
  \Delta L
  \;=\;
  \Delta L_\Sp + \Delta L_\Ge + \Delta L_\To + \Delta L_\Al ,
\]
where each $\Delta L_\lambda$ is the contribution of layer
$\lambda \in \{\Sp, \Ge, \To, \Al\}$ to the discrepancy between
$L^\ast$ and $L_\Cl$.
\end{definition}

\begin{proposition}[Layerwise correction]\label{prop:layerwise}
If each layer $\lambda$ is information-preserving in the sense of
Part~I, Definition~1.9, then
\[
  \Delta L_\lambda = 0
  \qquad \forall \lambda \in \{\Sp, \Ge, \To, \Al\},
\]
and hence $L^\ast = L_\Cl$.
\end{proposition}

\begin{proof}
Information preservation at layer $\lambda$ implies that no
structural data is lost in the transition $\phi_{\lambda, \lambda'}$.
Therefore the reconstruction $\rho_\Inv$ recovers $L_\Cl$ exactly,
and $\Delta L = 0$.
\end{proof}

% ------------------------------------------------------------
\section{Correction Operators and Correspondence}
\label{sec:correction}
% ------------------------------------------------------------

\subsection{The Correspondence Principle}

\begin{axiom}[Correspondence principle]\label{ax:correspondence}
The corrected operator $L_\Corr$ depends on a parameter
$\epsilon > 0$ and satisfies
\[
  \lim_{\epsilon \to 0^+} L_\Corr(\epsilon) = L_\Cl
\]
in the topology $\tau$ declared in Definition~\ref{def:classical-layer}.
\end{axiom}

\begin{remark}
Axiom~\ref{ax:correspondence} is the analogue, in the reverse
direction, of Axiom~1.5 of Part~I. It expresses the requirement that
the correction vanishes in the regime where the classical model is
valid.
\end{remark}

\subsection{Perturbative Regime}

\begin{definition}[Perturbative regime]
The correction is said to be \emph{perturbative} if
\[
  \Norm{\Delta L(\epsilon)} = \mathcal{O}(\epsilon)
  \qquad \text{as } \epsilon \to 0^+ .
\]
\end{definition}

\begin{definition}[Non-perturbative regime]
The correction is said to be \emph{non-perturbative} if
\[
  \Norm{\Delta L(\epsilon)} \not\to 0
  \qquad \text{as } \epsilon \to 0^+ .
\]
\end{definition}

\begin{remark}
Non-perturbative corrections are the primary objects of interest:
they signal structural discrepancies that cannot be absorbed by
smooth deformation of the classical model. In the terminology of
Part~I, they correspond to entries of the validation matrix marked
$\mathrm{O}$ in the reconstruction row.
\end{remark}

\subsection{Correspondence in the Spectral Layer}

\begin{proposition}[Spectral correspondence]\label{prop:spec-corr}
Suppose the correction $\Delta L$ is self-adjoint. Then
\[
  \spec(L_\Corr(\epsilon))
  \;\longrightarrow\;
  \spec(L_\Cl)
  \qquad \text{(Hausdorff convergence)}
\]
as $\epsilon \to 0^+$.
\end{proposition}

\begin{proof}
Apply Weyl's perturbation bound (Part~II, Theorem~2.10) to the
decomposition $L_\Corr(\epsilon) = L_\Cl + \Delta L(\epsilon)$ and
use $\Norm{\Delta L(\epsilon)} \to 0$.
\end{proof}

\subsection{Correspondence in the Geometric Layer}

\begin{definition}[Metric correction]
A metric correction is a family of Riemannian metrics
$\{g_\epsilon\}_{\epsilon > 0}$ on a fixed manifold $M$ such that
\[
  g_\epsilon \to g_\Cl
  \qquad \text{in } C^\infty_{\mathrm{loc}}(M, \mathrm{Sym}^2 T^\ast M) .
\]
\end{definition}

\begin{proposition}[Curvature stability]
If $g_\epsilon \to g_\Cl$ in $C^\infty_{\mathrm{loc}}$ then
\[
  R_\epsilon \to R_\Cl, \qquad
  \op{Ric}_\epsilon \to \op{Ric}_\Cl, \qquad
  S_\epsilon \to S_\Cl
\]
in $C^\infty_{\mathrm{loc}}$.
\end{proposition}

\begin{proof}
The curvature tensor is a polynomial in the first and second
derivatives of the metric. Local uniform convergence of $g_\epsilon$
and its derivatives to all orders therefore implies convergence of
$R_\epsilon$, $\op{Ric}_\epsilon$, $S_\epsilon$.
\end{proof}

% ------------------------------------------------------------
\section{Bidirectional Characterisation}
\label{sec:bidirectional}
% ------------------------------------------------------------

\subsection{The Bidirectional Diagram}

\begin{definition}[Bidirectional characterisation]
The methodology yields a \emph{bidirectional characterisation} if
there exist functors
\[
  \mathbf{C} : \Ocal_\Sp \to \Ocal_\Fo, \qquad
  \mathbf{R} : \Ocal_\Fo \to \Ocal_\Sp
\]
such that
\[
  \mathbf{R} \circ \mathbf{C} \;\simeq\; \mathrm{id}_{\Ocal_\Sp}, \qquad
  \mathbf{C} \circ \mathbf{R} \;\simeq\; \mathrm{id}_{\Ocal_\Fo},
\]
in the sense of asymptotic equivalence (Part~I, Definition~1.15).
\end{definition}

\begin{theorem}[Bidirectional characterisation of the correction]\label{thm:bidir}
Suppose:
\begin{enumerate}[label=(\roman*)]
  \item the construction functor $\mathbf{C}$ of Part~II,
        \S\ref{sec:composition}, is coherent;
  \item the inverse operator $\rho_\Inv$ is injective on
        $\Ecal_\Cl$;
  \item the correspondence principle (Axiom~\ref{ax:correspondence})
        holds;
  \item the reconstruction is stable in the sense of
        \S\ref{sec:wellposedness}.
\end{enumerate}
Then there exists a bidirectional characterisation
\[
  \Ob \;\rightleftarrows\; \Ecal_\Corr ,
\]
with the classical theory $\Ecal_\Cl$ appearing as the limit
$\epsilon \to 0^+$.
\end{theorem}

\begin{proof}[Sketch of proof]
Injectivity of $\rho_\Inv$ gives a well-defined inverse
$\rho_\Inv^{-1}$ on its image. The correspondence principle ensures
that the composite
\[
  \mathbf{R} \circ \mathbf{C} : \Ocal_\Sp \to \Ecal_\Corr \to \Ocal_\Sp
\]
is well-defined and, by coherence, asymptotic to the identity. The
reverse composite is handled symmetrically. The limit statement is
Axiom~\ref{ax:correspondence}.
\end{proof}

\subsection{Failure Modes}

\begin{definition}[Failure modes of the bidirectional characterisation]
The bidirectional characterisation fails if one of the following
occurs:
\begin{enumerate}[label=(\alph*)]
  \item \textbf{Loss of injectivity:} $\rho_\Inv$ is not injective,
        so distinct structures map to the same corrected model;
  \item \textbf{Breakdown of coherence:} some square of the
        bidirectional diagram fails to commute beyond the tolerance
        $\varepsilon_N$;
  \item \textbf{Non-vanishing correction:} $\Delta L(\epsilon)$ does
        not tend to zero as $\epsilon \to 0^+$;
  \item \textbf{Instability:} small perturbations of the input
        produce large perturbations of the output (see
        \S\ref{sec:wellposedness}).
\end{enumerate}
\end{definition}

\begin{proposition}[Diagnostic value of failure]
Each failure mode (a)--(d) identifies a specific deficiency in the
methodology, located at a specific layer and recorded in the
validation matrix $\Vmat$.
\end{proposition}

\subsection{Reconstruction and the Validation Matrix}

\begin{definition}[Reconstruction rows]
The reconstruction pipeline contributes three rows to the validation
matrix $\Vmat$, corresponding to the layers
\[
  \Cl, \quad \Corr, \quad \An .
\]
\end{definition}

\begin{center}
\begin{tabular}{c|c|l}
  Layer & Criterion & Status \\
  \hline
  $\Cl$   & Existence   & P \\
  $\Cl$   & Coherence   & P \\
  $\Cl$   & Invariance  & P \\
  $\Cl$   & Stability   & O \\
  $\Corr$ & Existence   & O \\
  $\Corr$ & Coherence   & O \\
  $\Corr$ & Correspondence & P (by axiom) \\
  $\An$   & Regularity  & O \\
  $\An$   & Analyticity & O \\
\end{tabular}
\end{center}

% ------------------------------------------------------------
\section{Well-Posedness, Stability, and Regularisation}
\label{sec:wellposedness}
% ------------------------------------------------------------

\subsection{Hadamard Well-Posedness}

\begin{definition}[Hadamard well-posedness]
The structural inverse problem (Definition~\ref{def:inverse-problem})
is \emph{well-posed} in the sense of Hadamard if:
\begin{enumerate}[label=(\roman*)]
  \item a solution exists for every classical model
        $\Mcal_\star \in \Ecal_\Cl$;
  \item the solution is unique;
  \item the solution depends continuously on the data.
\end{enumerate}
If any of these conditions fails, the problem is \emph{ill-posed}.
\end{definition}

\begin{proposition}[Failure of well-posedness]\label{prop:illposed}
The structural inverse problem is ill-posed whenever
$\rho_\Inv$ is not injective or not continuous.
\end{proposition}

\begin{proof}
Failure of injectivity violates condition (ii); failure of
continuity violates condition (iii).
\end{proof}

\subsection{Stability Estimates}

\begin{definition}[Stability estimate]
A \emph{stability estimate} for $\rho_\Inv$ is an inequality
\[
  d_{\Sp}(\rho_\Inv(\Mcal_1), \rho_\Inv(\Mcal_2))
  \;\leq\;
  \omega\bigl(d_{\Cl}(\Mcal_1, \Mcal_2)\bigr),
\]
where $\omega : [0, \infty) \to [0, \infty)$ is a modulus of
continuity with $\omega(0) = 0$.
\end{definition}

\begin{definition}[H\"older stability]
If $\omega(t) = C t^\alpha$ for some $C > 0$ and $\alpha \in (0, 1]$,
the stability is \emph{H\"older} of exponent $\alpha$.
\end{definition}

\begin{definition}[Logarithmic stability]
If $\omega(t) = C \Abs{\log t}^{-\beta}$ for some $C, \beta > 0$,
the stability is \emph{logarithmic} of exponent $\beta$.
\end{definition}

\begin{remark}
Logarithmic stability is the typical regime for severely ill-posed
inverse problems. H\"older stability is the typical regime for
moderately ill-posed problems. Both are weaker than Lipschitz
stability, which corresponds to $\alpha = 1$ and $C < \infty$.
\end{remark}

\subsection{Regularisation}

\begin{definition}[Regularisation operator]
A \emph{regularisation} of $\rho_\Inv$ is a family of operators
\[
  R_\alpha : \Ecal_\Cl \to \Ecal_\Corr, \qquad \alpha > 0,
\]
such that:
\begin{enumerate}[label=(\roman*)]
  \item $R_\alpha$ is continuous for every $\alpha > 0$;
  \item $R_\alpha \to \rho_\Inv$ as $\alpha \to 0^+$ on the domain
        of $\rho_\Inv$;
  \item the convergence in (ii) is uniform on compact subsets of the
        domain.
\end{enumerate}
\end{definition}

\begin{theorem}[Regularisation and stability]\label{thm:regularisation}
Let $R_\alpha$ be a regularisation of $\rho_\Inv$. Then for every
classical model $\Mcal_\star$ and every $\delta > 0$ there exists
$\alpha(\delta) > 0$ such that
\[
  \Norm{ R_{\alpha(\delta)}(\Mcal_\star) - \rho_\Inv(\Mcal_\star) }
  \;\leq\; \delta .
\]
\end{theorem}

\begin{proof}
By property (ii) of Definition~\ref{def:correction}--analogue for
$R_\alpha$, for each fixed $\Mcal_\star$ the quantity
$\Norm{ R_\alpha(\Mcal_\star) - \rho_\Inv(\Mcal_\star) }$ tends to
$0$ as $\alpha \to 0^+$. Choose $\alpha(\delta)$ accordingly.
\end{proof}

\begin{remark}
The choice $\alpha = \alpha(\delta)$ is the central trade-off of
regularisation theory: larger $\alpha$ gives greater stability but
greater bias; smaller $\alpha$ gives smaller bias but greater
variance. The optimal choice depends on the noise level $\delta$ and
the a priori smoothness of the solution.
\end{remark}

\subsection{Discretisation and Stability}

\begin{definition}[Discrete regularisation]
A \emph{discrete regularisation} of $\rho_\Inv$ is a family
\[
  R_{\alpha, N} : \Ecal_\Cl \to \Ecal_\Corr, \qquad \alpha > 0, \; N \in \mathbb{N},
\]
satisfying:
\begin{enumerate}[label=(\roman*)]
  \item $R_{\alpha, N}$ is continuous and finite-dimensional for each
        $\alpha, N$;
  \item $R_{\alpha, N} \to R_\alpha$ as $N \to \infty$;
  \item the discretisation error $\varepsilon_N$ satisfies
        $\varepsilon_N \to 0$ as $N \to \infty$.
\end{enumerate}
\end{definition}

\begin{proposition}[Consistency--stability--convergence for inversion]
Let $R_{\alpha, N}$ be a discrete regularisation. Then
\[
  \text{consistency} + \text{stability}
  \;\Longrightarrow\;
  \text{convergence}
\]
in the sense that
\[
  \lim_{N \to \infty} R_{\alpha(\delta), N}(\Mcal_\star)
  = \rho_\Inv(\Mcal_\star) .
\]
\end{proposition}

\begin{proof}
Standard argument in the spirit of Lax--Richtmyer
(Part~II, Theorem~\ref{thm:lax-rich}--analogue for inversion):
consistency controls the truncation error, stability controls the
propagation of perturbations, and their combination yields
convergence.
\end{proof}

\subsection{Analytic Verification}

\begin{definition}[Analytic layer]
The \emph{analytic layer} is the pair
\[
  \Acal = (\mathcal{A}(\Omega), \mathcal{D}),
\]
where $\mathcal{A}(\Omega)$ is a space of analytic functions on
$\Omega$ and $\mathcal{D}$ is a differential structure on
$\mathcal{A}(\Omega)$.
\end{definition}

\begin{definition}[Analyticity of the corrected solution]
The corrected solution $u_\Corr$ is \emph{analytic} if
$u_\Corr \in \mathcal{A}(\Omega)$ and
\[
  L_\Corr u_\Corr = 0
\]
holds in the sense of the differential structure $\mathcal{D}$.
\end{definition}

\begin{theorem}[Analytic closure]\label{thm:analytic-closure}
If
\begin{enumerate}[label=(\roman*)]
  \item $L_\Corr$ has analytic coefficients on $\Omega$;
  \item $u_\Corr$ satisfies $L_\Corr u_\Corr = 0$ weakly on $\Omega$;
  \item $L_\Corr$ is elliptic;
\end{enumerate}
then $u_\Corr$ is analytic on $\Omega$.
\end{theorem}

\begin{proof}
Apply the analytic regularity theorem for elliptic operators (see
Morrey, \emph{Multiple Integrals in the Calculus of Variations},
Chapter~6; or Bers--John--Schechter, \emph{Partial Differential
Equations}, Chapter~5).
\end{proof}

% ------------------------------------------------------------
\section{Composition of the Full Methodology}
\label{sec:composition-III}
% ------------------------------------------------------------

\subsection{The Full Diagram}

The full multirepresentational methodology assembles into the
following diagram.

\begin{equation}\label{eq:full-diagram}
\begin{tikzcd}[row sep=large, column sep=large]
  & \Ocal_\Sp \arrow[dl, "\phi_{\Sp\Ge}"'] \arrow[dr, "\mathbf{C}"'] & \\
  \Ocal_\Ge \arrow[d, "\phi_{\Ge\To}"'] & & \Ocal_\Fo \arrow[dll, "\rho_\Cl"'] \\
  \Ocal_\To \arrow[d, "\phi_{\To\Al}"'] & & \Ecal_\Cl \arrow[d, "\rho_\Inv"'] \\
  \Ocal_\Al \arrow[d, "\phi_{\Al\Nu}"'] & & \Ecal_\Corr \arrow[d, "\rho_\An"'] \\
  \Ocal_\Nu \arrow[dr, "\phi_{\Nu\Fo}"'] & & \Acal
\end{tikzcd}
\end{equation}

\subsection{Composite Operators}

\begin{definition}[Full construction--reconstruction cycle]
The \emph{full cycle} is the composite
\[
  \mathbf{K}
  \;=\;
  \rho_\An \circ \rho_\Inv \circ \rho_\Cl \circ \mathbf{C}
  :
  \Ocal_\Sp \to \Acal .
\]
\end{definition}

\begin{proposition}[Factorisation]\label{prop:factorisation}
The cycle $\mathbf{K}$ factors as
\[
  \mathbf{K} = \mathbf{K}_\An \circ \mathbf{K}_\Corr \circ \mathbf{K}_\Cl ,
\]
where
\[
  \mathbf{K}_\Cl   = \rho_\Cl \circ \mathbf{C}, \qquad
  \mathbf{K}_\Corr = \rho_\Inv, \qquad
  \mathbf{K}_\An   = \rho_\An .
\]
\end{proposition}

\subsection{Closure Condition}

\begin{definition}[Closure of the cycle]
The cycle $\mathbf{K}$ is \emph{closed} if there exists a natural
transformation
\[
  \eta : \mathrm{id}_{\Ocal_\Sp} \Longrightarrow \mathbf{K}
\]
or, more strongly, an asymptotic equivalence
\[
  \mathbf{K} \;\sim_\infty\; \mathrm{id}_{\Ocal_\Sp} .
\]
\end{definition}

\begin{theorem}[Closure criterion]\label{thm:closure-criterion}
The cycle $\mathbf{K}$ is closed if and only if:
\begin{enumerate}[label=(\roman*)]
  \item the correction $\Delta L(\epsilon)$ vanishes as
        $\epsilon \to 0^+$;
  \item the reconstruction $\rho_\Inv$ is injective;
  \item the correspondence between $\Ocal_\Sp$ and $\Acal$ is
        compatible with the analytic structure on both sides.
\end{enumerate}
\end{theorem}

\begin{proof}
(i) is necessary by Axiom~\ref{ax:correspondence}. (ii) is necessary
by Theorem~\ref{thm:bidir}. (iii) is necessary by
Theorem~\ref{thm:analytic-closure}. Sufficiency follows by composing
the three implications.
\end{proof}

\begin{remark}
The closure criterion is the ultimate requirement of the
methodology. When it holds, the classical theory $\Ecal_\Cl$
appears as the $\epsilon \to 0^+$ limit, and the corrected theory
$\Ecal_\Corr$ is a \emph{genuine extension}: it preserves the
classical limit while adding structural content not present in the
classical model.
\end{remark}

% ------------------------------------------------------------
\section*{Summary of Part III}
% ------------------------------------------------------------

Part~III has developed the reconstruction pipeline:
\begin{itemize}
  \item extraction of the classical model
        $\Ecal_\Cl = \rho_\Cl(\Ocal_\Fo)$
        (Definition~\ref{def:classical-layer});
  \item formulation of the structural inverse problem
        (Definition~\ref{def:inverse-problem});
  \item construction of the delta operator $\Delta L$ and the
        corrected operator $L_\Corr$ (Definition~\ref{def:correction});
  \item proof of the bidirectional characterisation
        $\Ob \rightleftarrows \Ecal_\Corr$
        (Theorem~\ref{thm:bidir});
  \item treatment of well-posedness, stability, and regularisation
        (\S\ref{sec:wellposedness});
  \item closure of the full construction--reconstruction cycle
        (Theorem~\ref{thm:closure-criterion}).
\end{itemize}

The key structural results are:
\[
  \mathcal{E}_\Cl \;\trianglelefteq\; \Ecal_\Cl
  \;\Longrightarrow\;
  \Delta L = 0,
\]
\[
  \rho_\Inv \text{ injective}
  \;\Longleftrightarrow\;
  \Ob \rightleftarrows \Ecal_\Corr \text{ is well-defined},
\]
and
\[
  \text{consistency} + \text{stability}
  \;\Longrightarrow\;
  \text{convergence}
\]
for both the direct and the inverse problems.

Part~IV will treat the discrete-to-continuum passage, the analytic
verification of the corrected model, and the validation matrix for
the full methodology.

\documentclass[11pt,a4paper]{article}

% ============================================================
% PACKAGES
% ============================================================
\usepackage[utf8]{inputenc}
\usepackage[T1]{fontenc}
\usepackage[english]{babel}
\usepackage{amsmath,amssymb,amsthm,mathtools}
\usepackage{mathrsfs}
\usepackage{tikz-cd}
\usepackage{enumitem}
\usepackage{hyperref}
\usepackage{geometry}
\usepackage{microtype}

\geometry{margin=1in}

% ============================================================
% THEOREM ENVIRONMENTS
% ============================================================
\theoremstyle{definition}
\newtheorem{definition}{Definition}[section]
\newtheorem{axiom}[definition]{Axiom}
\newtheorem{example}[definition]{Example}
\newtheorem{remark}[definition]{Remark}
\newtheorem{convention}[definition]{Convention}
\newtheorem{construction}[definition]{Construction}
\newtheorem{assumption}[definition]{Assumption}

\theoremstyle{plain}
\newtheorem{proposition}[definition]{Proposition}
\newtheorem{theorem}[definition]{Theorem}
\newtheorem{lemma}[definition]{Lemma}
\newtheorem{corollary}[definition]{Corollary}

% ============================================================
% MACROS
% ============================================================
\newcommand{\Lam}{\Lambda}
\newcommand{\Sp}{\mathrm{Sp}}
\newcommand{\Ge}{\mathrm{Ge}}
\newcommand{\To}{\mathrm{To}}
\newcommand{\Al}{\mathrm{Al}}
\newcommand{\Nu}{\mathrm{Nu}}
\newcommand{\Fo}{\mathrm{Fo}}
\newcommand{\Cl}{\mathrm{Cl}}
\newcommand{\Corr}{\mathrm{Corr}}
\newcommand{\Inv}{\mathrm{Inv}}
\newcommand{\An}{\mathrm{An}}
\newcommand{\Ob}{\mathfrak{O}}
\newcommand{\Ecal}{\mathcal{E}}
\newcommand{\Ocal}{\mathcal{O}}
\newcommand{\Acal}{\mathcal{A}}
\newcommand{\Mcal}{\mathcal{M}}
\newcommand{\Hcal}{\mathcal{H}}
\newcommand{\Dcal}{\mathcal{D}}
\newcommand{\Ncal}{\mathcal{N}}
\newcommand{\Ccal}{\mathcal{C}}
\newcommand{\philam}{\phi_{\lambda\mu}}
\newcommand{\rholam}{\rho_{\lambda}}
\newcommand{\Vmat}{\mathbf{V}}
\newcommand{\Jset}{\mathcal{J}}
\newcommand{\dd}{\,\mathrm{d}}
\newcommand{\op}[1]{\operatorname{#1}}
\newcommand{\spec}{\sigma}
\newcommand{\hilb}{\mathcal{H}}
\newcommand{\Cat}{\mathbf{Cat}}
\newcommand{\Set}{\mathbf{Set}}
\newcommand{\Vect}{\mathbf{Vect}}
\newcommand{\Grp}{\mathbf{Grp}}
\newcommand{\Rng}{\mathbf{Rng}}
\newcommand{\Top}{\mathbf{Top}}
\newcommand{\Man}{\mathbf{Man}}
\newcommand{\Norm}[1]{\left\lVert #1 \right\rVert}
\newcommand{\Abs}[1]{\left\lvert #1 \right\rvert}
\newcommand{\ip}[2]{\left\langle #1, #2 \right\rangle}
\newcommand{\R}{\mathbb{R}}
\newcommand{\C}{\mathbb{C}}
\newcommand{\N}{\mathbb{N}}
\newcommand{\Z}{\mathbb{Z}}

% ============================================================
% TITLE
% ============================================================
\title{%
  \textbf{A Multirepresentational Construction--Reconstruction\\
  Methodology for Mathematical Structures}\\[4pt]
  \large Part IV: Discrete-to-Continuum Passage and Analytic Verification
}
\author{}
\date{}

\begin{document}
\maketitle

\begin{abstract}
\noindent
Part~IV addresses the discrete-to-continuum passage of the
multirepresentational methodology. Starting from the discrete families
$\Ob^{(N)}$ introduced in Part~I, we develop a rigorous framework for
the study of structural emergence, stability, and convergence as the
resolution parameter $N \to \infty$. The central objects are the
layerwise convergence statements
$\Ocal_\lambda^{(N)} \to \Ocal_\lambda$, the emergent invariants
detected by numerical computation, the compactness and uniform-bound
arguments required to pass to the continuum limit, and the final
regularity and analyticity statements that close the methodology.
The chapter distinguishes sharply between numerical evidence and
mathematical proof, and records the passage from the former to the
latter as a sequence of four steps: detection, estimation,
compactness, and limit identification.
\end{abstract}

\tableofcontents

% ============================================================
\part{Discrete-to-Continuum Passage and Analytic Verification}
% ============================================================

\section{Overview of the Discrete-to-Continuum Programme}
\label{sec:overview-IV}

The construction pipeline of Part~II and the reconstruction pipeline
of Part~III are, so far, formal statements about the layers
$\Ocal_\lambda$. They do not, by themselves, produce computable or
verifiable content. That content is supplied by the \emph{discrete
families}
\[
  \Ob^{(N)} = \bigl(\Ocal_\lambda^{(N)}\bigr)_{\lambda \in \Lam},
  \qquad N \in \N,
\]
introduced in Part~I, \S1.6.

Part~IV is organised as follows.
Section~\ref{sec:discrete-families} defines discrete families and
their nested structure. Section~\ref{sec:stability} treats structural
stability and invariant detection. Section~\ref{sec:csc} develops the
consistency--stability--convergence triangle. Section~\ref{sec:compactness}
supplies the compactness and uniform-bound machinery required for the
continuum limit. Section~\ref{sec:regularity} closes with regularity
and analyticity statements.

\begin{remark}
The discrete-to-continuum programme is the layer at which the
methodology interacts with actual computation. It is therefore the
layer at which the distinction between \emph{evidence} and
\emph{proof} becomes critical, and it is the layer that gives the
validation matrix $\Vmat$ its off-diagonal entries marked
$\mathrm{N}$.
\end{remark}

% ------------------------------------------------------------
\section{Discrete Families}
\label{sec:discrete-families}
% ------------------------------------------------------------

\subsection{The Resolution Parameter}

\begin{definition}[Resolution parameter]\label{def:resolution}
A \emph{resolution parameter} is a positive integer $N \in \N$
indexing a family of finite-dimensional approximations of the master
object $\Ob$. The integer $N$ is assumed to increase monotonically
along a filtration
\[
  \Ncal_1 \subset \Ncal_2 \subset \cdots \subset \Ncal_N \subset \cdots
\]
of nested finite-dimensional subspaces of the ambient Hilbert space
$\hilb$.
\end{definition}

\begin{definition}[Discrete master object]\label{def:discrete-master}
The \emph{discrete master object} at resolution $N$ is
\[
  \Ob^{(N)}
  \;=\;
  \bigl(\Ocal_\Sp^{(N)}, \Ocal_\Ge^{(N)}, \Ocal_\To^{(N)},
        \Ocal_\Al^{(N)}, \Ocal_\Nu^{(N)}, \Ocal_\Fo^{(N)}\bigr),
\]
together with layer projections
\[
  \pi_\lambda^{(N)} : \Ob^{(N)} \to \Ocal_\lambda^{(N)}
\]
and resolution morphisms
\[
  \iota_N : \Ob^{(N)} \to \Ob .
\]
\end{definition}

\subsection{Nested Approximations}

\begin{axiom}[Nestedness]\label{ax:nested}
For every $N \leq M$ there exist canonical inclusions
\[
  \iota_{N,M} : \Ocal_\lambda^{(N)} \hookrightarrow \Ocal_\lambda^{(M)},
  \qquad \lambda \in \Lam,
\]
compatible with the layer projections:
\[
  \pi_\lambda^{(M)} \circ \iota_{N,M}
  \;=\;
  \iota_{N,M} \circ \pi_\lambda^{(N)} .
\]
\end{axiom}

\begin{definition}[Inductive limit]
The \emph{inductive limit} of the discrete family is
\[
  \Ocal_\lambda^{(\infty)}
  \;=\;
  \varinjlim_{N} \Ocal_\lambda^{(N)} .
\]
\end{definition}

\begin{proposition}[Universal property]\label{prop:universal}
The inductive limit $\Ocal_\lambda^{(\infty)}$ is the unique
object, up to isomorphism, together with canonical morphisms
$\jmath_N : \Ocal_\lambda^{(N)} \to \Ocal_\lambda^{(\infty)}$
satisfying
\[
  \jmath_M \circ \iota_{N,M} = \jmath_N
  \qquad (N \leq M),
\]
and universal with respect to this property.
\end{proposition}

\begin{proof}
Standard category-theoretic construction of the colimit; existence
follows from nestedness (Axiom~\ref{ax:nested}) and uniqueness from
the universal property.
\end{proof}

\subsection{Emergent Structures}

\begin{definition}[Emergent structure]\label{def:emergent}
A structure $\mathcal{P}$ on the limit object
$\Ocal_\lambda^{(\infty)}$ is \emph{emergent} from the discrete
family $\{\Ocal_\lambda^{(N)}\}_N$ if:
\begin{enumerate}[label=(\roman*)]
  \item there exist structures $\mathcal{P}_N$ on each
        $\Ocal_\lambda^{(N)}$ such that
        $\mathcal{P}_N \to \mathcal{P}$ in an appropriate topology;
  \item the convergence is stable under the inclusions
        $\iota_{N,M}$.
\end{enumerate}
\end{definition}

\begin{definition}[Persistent structure]\label{def:persistent}
A structure $\mathcal{P}$ is \emph{persistent} at level $N_0$ if
$\mathcal{P}_N \neq \varnothing$ for every $N \geq N_0$.
\end{definition}

\begin{example}[Persistent homology]
In topological data analysis, the persistent homology of a filtration
\[
  X_1 \subset X_2 \subset \cdots \subset X_N
\]
produces a barcode
\[
  \mathrm{PH}_k(X_\bullet) = \{(b_i, d_i)\}_i
\]
of birth--death intervals. Structures whose intervals remain finite
as $N \to \infty$ are the topological analogues of emergent
structures in the sense of Definition~\ref{def:emergent}.
\end{example}

\begin{remark}
The notion of emergent structure is precisely the mathematical
formalisation of the informal notion of ``pattern'': a pattern is a
structure that persists across resolutions and admits a well-defined
limit.
\end{remark}

% ------------------------------------------------------------
\section{Stability and Invariance}
\label{sec:stability}
% ------------------------------------------------------------

\subsection{Structural Stability}

\begin{definition}[Structural stability]\label{def:structural-stability}
Let $\Ocal_\lambda^{(N)}$ and $\Ocal_\lambda^{(N)\prime}$ be two
discrete approximations at resolution $N$ arising from perturbations
of the input data. The layer $\lambda$ is \emph{structurally stable}
if for every $\epsilon > 0$ there exists $\delta > 0$ such that
\[
  d_{\mathrm{in}}\bigl(\mathrm{data}, \mathrm{data}'\bigr) < \delta
  \;\Longrightarrow\;
  d_\lambda\bigl(\Ocal_\lambda^{(N)}, \Ocal_\lambda^{(N)\prime}\bigr) < \epsilon
\]
uniformly in $N$.
\end{definition}

\begin{definition}[Uniform structural stability]
The family $\{\Ocal_\lambda^{(N)}\}_N$ is \emph{uniformly
structurally stable} if the modulus of continuity is independent
of $N$.
\end{definition}

\subsection{Robustness}

\begin{definition}[Robust invariant]\label{def:robust}
An invariant $I_\lambda : \Ocal_\lambda \to \mathcal{V}_\lambda$ is
\emph{robust} if it is continuous with respect to the layer topology
and satisfies
\[
  I_\lambda\bigl(\Ocal_\lambda^{(N)}\bigr)
  \;\longrightarrow\;
  I_\lambda\bigl(\Ocal_\lambda^{(\infty)}\bigr)
\]
uniformly in the admissible class of inputs.
\end{definition}

\begin{proposition}[Robustness and emergence]\label{prop:robust-emergent}
If $I_\lambda$ is robust then the structure
$\mathcal{P} = I_\lambda^{-1}(v)$ for fixed $v \in \mathcal{V}_\lambda$
is emergent in the sense of Definition~\ref{def:emergent}.
\end{proposition}

\begin{proof}
Robustness gives the convergence $I_\lambda(\Ocal_\lambda^{(N)}) \to v$
along the filtration. Emergence follows by choosing
$\mathcal{P}_N = I_\lambda^{-1}(v) \cap \Ocal_\lambda^{(N)}$.
\end{proof}

\subsection{Invariant Detection}

\begin{definition}[Detection algorithm]\label{def:detection}
An \emph{invariant detection algorithm} is a map
\[
  \Dcal : \{\Ocal_\lambda^{(N)}\}_{N=1}^{N_{\max}}
  \;\longrightarrow\;
  \mathcal{H},
\]
where $\mathcal{H}$ is a hypothesis space, that outputs a candidate
invariant $I_\lambda \in \mathcal{H}$ consistent with the discrete
data.
\end{definition}

\begin{remark}
Invariant detection is the formalisation of the informal notion of
``spotting a pattern''. The output $I_\lambda$ of a detection
algorithm is a \emph{conjecture}, not a theorem: it must subsequently
be verified by an estimation argument (\S\ref{sec:csc}) and a
compactness argument (\S\ref{sec:compactness}).
\end{remark}

\begin{example}[Numerical detection of $\pi$]
Consider the sequence of approximations
\[
  I_N = 4 \sum_{k=0}^{N} \frac{(-1)^k}{2k+1} .
\]
A detection algorithm may observe $I_N \to 3.14159\ldots$ and
hypothesise $I = \pi$. The conjecture is verified by proving
\[
  \Abs{I_N - \pi} \leq \frac{C}{N},
\]
which is a classical estimate.
\end{example}

\begin{remark}
The detection algorithm outputs a plausible invariant. The
verification is what makes it a theorem. The distinction is
reflected in the validation matrix by the presence of both
$\mathrm{N}$ (numerical evidence) and $\mathrm{P}$ (proof) entries.
\end{remark}

% ------------------------------------------------------------
\section{Consistency, Stability, and Convergence}
\label{sec:csc}
% ------------------------------------------------------------

\subsection{Consistency}

\begin{definition}[Layerwise consistency]\label{def:consistency}
Let $\Ocal_\lambda$ be a layer and let $\Ocal_\lambda^{(N)}$ be its
$N$-th discrete approximation. Let
$\mathcal{F}_\lambda : \Ocal_\lambda \to \mathcal{W}_\lambda$
be the defining equation of the layer. The discretisation is
\emph{consistent} if
\[
  \Norm{
    \mathcal{F}_\lambda\bigl(\pi_\lambda^{(N)} u\bigr)
    -
    \pi_\lambda^{(N)} \mathcal{F}_\lambda(u)
  }_{\mathcal{W}_\lambda}
  \;\longrightarrow\;
  0
  \qquad (N \to \infty),
\]
uniformly on bounded subsets of the domain.
\end{definition}

\begin{definition}[Order of consistency]
If the convergence in Definition~\ref{def:consistency} is
$\mathcal{O}(N^{-p})$ for some $p > 0$, the discretisation is said
to be \emph{consistent of order $p$}.
\end{definition}

\subsection{Stability}

\begin{definition}[Layerwise stability]\label{def:stability}
The discretisation is \emph{stable} if there exists a constant
$C > 0$, independent of $N$, such that
\[
  \Norm{(\mathcal{F}_\lambda^{(N)})^{-1}}_{\mathcal{W}_\lambda \to \Ocal_\lambda^{(N)}}
  \;\leq\;
  C
  \qquad \forall N \in \N .
\]
\end{definition}

\begin{definition}[Condition number]
The \emph{condition number} of $\mathcal{F}_\lambda^{(N)}$ is
\[
  \kappa\bigl(\mathcal{F}_\lambda^{(N)}\bigr)
  \;=\;
  \Norm{\mathcal{F}_\lambda^{(N)}} \cdot
  \Norm{(\mathcal{F}_\lambda^{(N)})^{-1}} .
\]
\end{definition}

\begin{remark}
Stability is equivalent to uniform boundedness of the condition
number. Numerical instability manifests as $\kappa \to \infty$ as
$N \to \infty$, which signals a loss of well-posedness in the
discrete problem.
\end{remark}

\subsection{Convergence}

\begin{theorem}[Consistency--stability--convergence]\label{thm:csc}
Let $\Ocal_\lambda^{(N)}$ be a consistent and stable discretisation
of $\Ocal_\lambda$. Then for every $u \in \Ocal_\lambda$,
\[
  \Norm{\pi_\lambda^{(N)} u - u_N}_{\Ocal_\lambda^{(N)}}
  \;\leq\;
  C \, \Norm{
    \mathcal{F}_\lambda\bigl(\pi_\lambda^{(N)} u\bigr)
    -
    \pi_\lambda^{(N)} \mathcal{F}_\lambda(u)
  }_{\mathcal{W}_\lambda},
\]
where $u_N \in \Ocal_\lambda^{(N)}$ is the discrete solution. In
particular,
\[
  u_N \;\longrightarrow\; u
  \qquad (N \to \infty).
\]
\end{theorem}

\begin{proof}
Write the discretised equation as
\[
  \mathcal{F}_\lambda^{(N)} u_N = 0 .
\]
Subtract $\mathcal{F}_\lambda^{(N)} \pi_\lambda^{(N)} u$ from both
sides to obtain
\[
  \mathcal{F}_\lambda^{(N)}(u_N - \pi_\lambda^{(N)} u)
  =
  - \mathcal{F}_\lambda^{(N)} \pi_\lambda^{(N)} u .
\]
Apply the stability estimate (Definition~\ref{def:stability}) and
the consistency estimate (Definition~\ref{def:consistency}) to
conclude.
\end{proof}

\begin{corollary}[Order of convergence]\label{cor:order}
If the discretisation is consistent of order $p$ and stable, then
the convergence is of order $p$:
\[
  \Norm{\pi_\lambda^{(N)} u - u_N}
  \;=\;
  \mathcal{O}(N^{-p}) .
\]
\end{corollary}

\begin{proof}
Immediate from Theorem~\ref{thm:csc} and Definition~\ref{def:consistency}.
\end{proof}

\subsection{The Three-Proof Structure}

\begin{definition}[Three-proof structure]\label{def:three-proof}
A layer $\lambda$ is said to satisfy the \emph{three-proof
structure} if:
\begin{enumerate}[label=(\roman*)]
  \item consistency (Definition~\ref{def:consistency}) holds;
  \item stability (Definition~\ref{def:stability}) holds;
  \item convergence (Theorem~\ref{thm:csc}) holds.
\end{enumerate}
\end{definition}

\begin{proposition}[Three-proof sufficiency]\label{prop:three-proof}
If a layer satisfies the three-proof structure, then
\[
  \Ocal_\lambda^{(N)} \to \Ocal_\lambda
\]
in the sense of Part~I, Definition~1.7.
\end{proposition}

\begin{proof}
Apply Theorem~\ref{thm:csc} to each defining equation of the layer
and combine the resulting estimates into the aggregate error
$\varepsilon_N$ of Part~I, \eqref{eq:error}.
\end{proof}

\begin{remark}
The three-proof structure replaces the informal statement ``the
numerical results look convergent'' by three explicit mathematical
obligations. This is the operational content of the distinction
between numerical evidence and mathematical proof.
\end{remark}

% ------------------------------------------------------------
\section{Compactness and Limit Passage}
\label{sec:compactness}
% ------------------------------------------------------------

\subsection{Uniform Bounds}

\begin{assumption}[Uniform boundedness]\label{ass:uniform}
There exists $C > 0$, independent of $N$, such that
\[
  \Norm{u_N}_{X}
  \;\leq\;
  C
  \qquad \forall N \in \N,
\]
where $X$ is a Banach space into which all $\Ocal_\lambda^{(N)}$
embed continuously.
\end{assumption}

\begin{proposition}[Consequence of uniform boundedness]\label{prop:uniform-bounded}
Under Assumption~\ref{ass:uniform}, the sequence $\{u_N\}_N$ admits
a weakly convergent subsequence in $X$:
\[
  u_{N_j} \rightharpoonup u_\infty \in X .
\]
\end{proposition}

\begin{proof}
Banach--Alaoglu theorem applied to the unit ball of $X$, using the
uniform bound to stay within a fixed ball.
\end{proof}

\subsection{Compactness Arguments}

\begin{theorem}[Compactness]\label{thm:compactness}
Let $X \hookrightarrow Y$ be a compact embedding of Banach spaces.
Under Assumption~\ref{ass:uniform}, there exists a subsequence
$\{u_{N_j}\}$ and a limit $u_\infty \in X$ such that
\[
  u_{N_j} \;\longrightarrow\; u_\infty
  \qquad \text{strongly in } Y .
\]
\end{theorem}

\begin{proof}
Weak compactness (Proposition~\ref{prop:uniform-bounded}) provides
a weakly convergent subsequence. Compactness of $X \hookrightarrow Y$
upgrades weak convergence in $X$ to strong convergence in $Y$.
\end{proof}

\begin{example}[Sobolev compactness]
The classical Rellich--Kondrachov theorem states that
\[
  H^1(\Omega) \hookrightarrow L^2(\Omega)
\]
is compact for $\Omega \subset \R^n$ bounded and Lipschitz. This
provides the compactness required in many applications of
Theorem~\ref{thm:compactness}.
\end{example}

\subsection{Limit Identification}

\begin{definition}[Limit equation]\label{def:limit-eq}
The \emph{limit equation} of the discrete family $\{u_N\}_N$ is
\[
  \mathcal{E}(u_\infty) = 0,
\]
where $\mathcal{E}$ is the formal equation obtained by passing to the
limit in the discretised equations $\mathcal{E}^{(N)}(u_N) = 0$.
\end{definition}

\begin{theorem}[Limit identification]\label{thm:limit-id}
Under the hypotheses of Theorem~\ref{thm:compactness}, the limit
$u_\infty$ satisfies the limit equation $\mathcal{E}(u_\infty) = 0$
in the weak sense of the space $Y$.
\end{theorem}

\begin{proof}
Use density of test functions in $Y^\ast$ and pass to the limit in
the weak formulation of $\mathcal{E}^{(N)}(u_N) = 0$. Consistency
(Definition~\ref{def:consistency}) ensures that the discrete operator
converges to the continuum operator weakly.
\end{proof}

\subsection{Extrapolation to the Limit}

\begin{definition}[Numerical extrapolation]\label{def:extrapolation}
Let $\{u_N\}_N$ be a convergent sequence with order $p$
(Corollary~\ref{cor:order}). The \emph{Richardson extrapolation} of
order $q$ is the linear combination
\[
  u_N^{(q)}
  \;=\;
  \sum_{j=0}^{q} c_j^{(q)} u_{N_j} ,
\]
where the coefficients $c_j^{(q)}$ are chosen so that the leading
$q$ error terms cancel.
\end{definition}

\begin{proposition}[Acceleration of convergence]\label{prop:richardson}
If $\{u_N\}$ converges with order $p$ and admits an asymptotic
expansion
\[
  u_N = u_\infty + a_1 N^{-p} + a_2 N^{-2p} + \cdots,
\]
then the Richardson extrapolation $u_N^{(q)}$ converges with order
$(q+1) p$.
\end{proposition}

\begin{proof}
By construction of the coefficients $c_j^{(q)}$, the first $q$ terms
of the asymptotic expansion cancel, leaving an error of order
$N^{-(q+1)p}$.
\end{proof}

\begin{remark}
Numerical extrapolation is a computational technique, not a
mathematical proof. It accelerates convergence on the discrete side
but does not, by itself, establish the limit statement. The latter
requires a mathematical estimate of the type
$\Norm{u_N - u_\infty} \leq C N^{-p}$, which is the content of
Corollary~\ref{cor:order}.
\end{remark}

% ------------------------------------------------------------
\section{Regularity and Analyticity}
\label{sec:regularity}
% ------------------------------------------------------------

\subsection{Regularity}

\begin{definition}[Sobolev regularity]\label{def:sobolev}
Let $\Omega \subset \R^n$ be an open set and let $s \in \R$. The
\emph{Sobolev space} $H^s(\Omega)$ is the completion of
$C_c^\infty(\Omega)$ with respect to the norm
\[
  \Norm{u}_{H^s}^2
  \;=\;
  \int_{\R^n} (1 + \Abs{\xi}^2)^s \Abs{\hat{u}(\xi)}^2 \dd \xi ,
\]
where $\hat{u}$ denotes the Fourier transform of $u$.
\end{definition}

\begin{definition}[Regularity of the limit]\label{def:regularity}
The limit $u_\infty$ is said to have \emph{regularity $s$} if
$u_\infty \in H^s(\Omega)$. It is said to have \emph{classical
regularity $C^k$} if $u_\infty \in C^k(\Omega)$ and all derivatives
up to order $k$ are continuous.
\end{definition}

\begin{theorem}[Regularity bootstrap]\label{thm:bootstrap}
Let $L$ be an elliptic operator of order $m$ with smooth
coefficients on $\Omega$, and let $u \in H^s(\Omega)$ satisfy
$Lu = f$ with $f \in H^{s-m}(\Omega)$. Then
\[
  u \in H^{s+1}(\Omega)
\]
whenever $s$ and $s+1$ are integers.
\end{theorem}

\begin{proof}
Standard elliptic regularity (see Agmon--Douglis--Nirenberg,
\emph{Estimates near the boundary for solutions of elliptic partial
differential equations}, Comm.~Pure Appl.~Math.~12 (1959),
623--727).
\end{proof}

\begin{corollary}[Smoothness under iteration]\label{cor:smoothness}
If the hypotheses of Theorem~\ref{thm:bootstrap} hold for every
$s \geq 0$, then $u \in C^\infty(\Omega)$.
\end{corollary}

\begin{proof}
Iterate the bootstrap: $u \in H^{s_0} \Rightarrow u \in H^{s_0+1}
\Rightarrow \cdots \Rightarrow u \in H^s$ for every $s$. Sobolev
embedding gives $u \in C^k$ for every $k$.
\end{proof}

\subsection{Analyticity}

\begin{definition}[Analytic function]\label{def:analytic}
A function $u : \Omega \to \C$ is \emph{analytic} if for every
$x_0 \in \Omega$ there exists $r > 0$ and coefficients
$\{a_\alpha\}_{\alpha \in \N^n}$ such that
\[
  u(x) = \sum_{\alpha} a_\alpha (x - x_0)^\alpha
  \qquad (\Abs{x - x_0} < r),
\]
where the series converges absolutely and uniformly on the ball of
radius $r$.
\end{definition}

\begin{theorem}[Analytic regularity for elliptic operators]\label{thm:analytic-regularity}
Let $L$ be an elliptic operator of order $m$ with real-analytic
coefficients on $\Omega$, and let $u \in H^s(\Omega)$ satisfy
$Lu = f$ with $f$ real-analytic on $\Omega$. Then $u$ is
real-analytic on $\Omega$.
\end{theorem}

\begin{proof}
See Morrey, \emph{Multiple Integrals in the Calculus of Variations},
Chapter~6, or Bers--John--Schechter, \emph{Partial Differential
Equations}, Chapter~5. The proof uses a bootstrapping argument
similar to that of Theorem~\ref{thm:bootstrap} but with analytic
norms replacing Sobolev norms.
\end{proof}

\subsection{Analytic Closure}

\begin{definition}[Analytic closure of the methodology]\label{def:analytic-closure}
The multirepresentational methodology is said to achieve
\emph{analytic closure} if the corrected solution $u_\Corr$
constructed in Part~III satisfies
\[
  u_\Corr \in \mathcal{A}(\Omega)
\]
and
\[
  L_\Corr u_\Corr = 0
\]
in the analytic sense of Definition~\ref{def:analytic}.
\end{definition}

\begin{theorem}[Analytic closure criterion]\label{thm:analytic-closure-IV}
Analytic closure holds if and only if:
\begin{enumerate}[label=(\roman*)]
  \item the corrected operator $L_\Corr$ is elliptic;
  \item its coefficients are real-analytic on $\Omega$;
  \item the discrete family $\{u_N\}_N$ is bounded in $H^s(\Omega)$
        for some $s > n/2$;
  \item the limit $u_\infty$ satisfies $L_\Corr u_\infty = 0$ in the
        weak sense.
\end{enumerate}
\end{theorem}

\begin{proof}
Necessity of (i)--(iv) follows from the definitions of elliptic
operator, analytic coefficient, uniform boundedness
(Assumption~\ref{ass:uniform}), and limit equation
(Definition~\ref{def:limit-eq}). Sufficiency follows from
Theorem~\ref{thm:analytic-regularity} applied to $L = L_\Corr$ and
$u = u_\infty$.
\end{proof}

\subsection{The Final Verification Chain}

\begin{definition}[Final verification chain]\label{def:final-chain}
The \emph{final verification chain} of the methodology is the
sequence
\[
  \text{detection}
  \;\to\;
  \text{estimation}
  \;\to\;
  \text{compactness}
  \;\to\;
  \text{limit}
  \;\to\;
  \text{regularity}
  \;\to\;
  \text{analyticity} .
\]
\end{definition}

\begin{theorem}[Closure of the verification chain]\label{thm:chain-closure}
If
\begin{enumerate}[label=(\roman*)]
  \item the detection algorithm outputs a robust invariant
        (Definition~\ref{def:robust});
  \item the three-proof structure holds (Definition~\ref{def:three-proof});
  \item the uniform bound holds (Assumption~\ref{ass:uniform});
  \item the compactness embedding is available (Theorem~\ref{thm:compactness});
  \item the analytic regularity hypotheses hold
        (Theorem~\ref{thm:analytic-regularity});
\end{enumerate}
then the full cycle
\[
  \Ob \;\longrightarrow\; \Ecal_\Corr \;\longrightarrow\; \Acal
\]
of Part~III terminates in a solution $u_\Corr \in \mathcal{A}(\Omega)$
satisfying $L_\Corr u_\Corr = 0$.
\end{theorem}

\begin{proof}
Combine the chain of implications given by
Propositions~\ref{prop:robust-emergent}, \ref{prop:three-proof};
Theorems~\ref{thm:compactness}, \ref{thm:limit-id}, and
\ref{thm:analytic-regularity}.
\end{proof}

\begin{remark}
Theorem~\ref{thm:chain-closure} is the closing statement of the
methodology. It asserts that the discrete-to-continuum passage,
combined with the analytic regularity theory, terminates in a
well-defined analytic solution of the corrected equation. When it
holds, the methodology is \emph{complete} in the sense that every
layer, from the spectral to the analytic, has been traversed and
verified.
\end{remark}

% ------------------------------------------------------------
\section{The Validation Matrix at the Discrete-to-Continuum Level}
\label{sec:validation-IV}
% ------------------------------------------------------------

\begin{definition}[Discrete-to-continuum rows]
The discrete-to-continuum programme contributes the following rows
to the validation matrix $\Vmat$, in addition to those already
recorded in Parts~I--III.
\end{definition}

\begin{center}
\begin{tabular}{c|c|l}
  Layer / stage & Criterion & Status \\
  \hline
  Detection            & Existence    & N \\
  Detection            & Invariance   & N \\
  Detection            & Coherence    & O \\
  Estimation           & Existence    & O \\
  Estimation           & Stability    & O \\
  Compactness          & Existence    & P (conditional) \\
  Limit passage        & Convergence  & P (conditional) \\
  Regularity           & Regularity   & P (conditional) \\
  Analyticity          & Regularity   & P (conditional) \\
  Analyticity          & Reversibility& O \\
\end{tabular}
\end{center}

\begin{remark}
The ``conditional'' entries indicate that the corresponding result
holds under the explicit hypotheses stated in the theorem. The
entry $\mathrm{N}$ in the detection row indicates that the
corresponding assertion is supported by numerical evidence but has
not been proved; the entry $\mathrm{O}$ indicates that the
assertion is open at the time of writing.
\end{remark}

\begin{proposition}[Necessary conditions for closure]\label{prop:necessary-closure}
If the validation matrix $\Vmat$ contains an entry marked
$\mathrm{F}$ (failed), then the methodology is not closed in the
sense of Part~I, Definition~1.19.
\end{proposition}

\begin{proof}
A failed entry indicates that some layer or stage of the
methodology does not satisfy its defining condition. Closure requires
every layer to satisfy at least one criterion marked $\mathrm{P}$
and no criterion marked $\mathrm{F}$; the presence of $\mathrm{F}$
violates the second requirement.
\end{proof}

% ------------------------------------------------------------
\section*{Summary of Part IV}
% ------------------------------------------------------------

Part~IV has developed the discrete-to-continuum passage of the
methodology:
\begin{itemize}
  \item discrete families $\Ob^{(N)}$ and the resolution parameter
        $N$ (Definition~\ref{def:resolution});
  \item nested approximations and inductive limits
        (Axiom~\ref{ax:nested}, Proposition~\ref{prop:universal});
  \item emergent and persistent structures
        (Definitions~\ref{def:emergent}, \ref{def:persistent});
  \item structural stability and robust invariants
        (Definitions~\ref{def:structural-stability}, \ref{def:robust});
  \item the consistency--stability--convergence triangle
        (Theorem~\ref{thm:csc});
  \item compactness and limit identification
        (Theorems~\ref{thm:compactness}, \ref{thm:limit-id});
  \item regularity and analyticity of the limit
        (Theorems~\ref{thm:bootstrap}, \ref{thm:analytic-regularity});
  \item the final verification chain
        (Definition~\ref{def:final-chain}) and its closure
        (Theorem~\ref{thm:chain-closure});
  \item the discrete-to-continuum rows of the validation matrix
        $\Vmat$.
\end{itemize}

The central structural results are:
\[
  \text{consistency} + \text{stability}
  \;\Longrightarrow\;
  \text{convergence},
\]
\[
  \text{uniform bound} + \text{compact embedding}
  \;\Longrightarrow\;
  \text{limit identification},
\]
and
\[
  \text{ellipticity} + \text{analytic coefficients}
  \;\Longrightarrow\;
  \text{analytic regularity}.
\]

Part~V will treat the validation matrix in full, the worked
example, and the formalisation of the methodology in proof
assistants.

\documentclass[11pt,a4paper]{article}

% ============================================================
% PACKAGES
% ============================================================
\usepackage[utf8]{inputenc}
\usepackage[T1]{fontenc}
\usepackage[english]{babel}
\usepackage{amsmath,amssymb,amsthm,mathtools}
\usepackage{mathrsfs}
\usepackage{tikz-cd}
\usepackage{enumitem}
\usepackage{hyperref}
\usepackage{geometry}
\usepackage{microtype}
\usepackage{booktabs}
\usepackage{array}
\usepackage{longtable}

\geometry{margin=1in}

% ============================================================
% THEOREM ENVIRONMENTS
% ============================================================
\theoremstyle{definition}
\newtheorem{definition}{Definition}[section]
\newtheorem{axiom}[definition]{Axiom}
\newtheorem{example}[definition]{Example}
\newtheorem{remark}[definition]{Remark}
\newtheorem{convention}[definition]{Convention}
\newtheorem{construction}[definition]{Construction}
\newtheorem{assumption}[definition]{Assumption}
\newtheorem{protocol}[definition]{Protocol}

\theoremstyle{plain}
\newtheorem{proposition}[definition]{Proposition}
\newtheorem{theorem}[definition]{Theorem}
\newtheorem{lemma}[definition]{Lemma}
\newtheorem{corollary}[definition]{Corollary}

% ============================================================
% MACROS
% ============================================================
\newcommand{\Lam}{\Lambda}
\newcommand{\Sp}{\mathrm{Sp}}
\newcommand{\Ge}{\mathrm{Ge}}
\newcommand{\To}{\mathrm{To}}
\newcommand{\Al}{\mathrm{Al}}
\newcommand{\Nu}{\mathrm{Nu}}
\newcommand{\Fo}{\mathrm{Fo}}
\newcommand{\Cl}{\mathrm{Cl}}
\newcommand{\Corr}{\mathrm{Corr}}
\newcommand{\Inv}{\mathrm{Inv}}
\newcommand{\An}{\mathrm{An}}
\newcommand{\Ob}{\mathfrak{O}}
\newcommand{\Ecal}{\mathcal{E}}
\newcommand{\Ocal}{\mathcal{O}}
\newcommand{\Acal}{\mathcal{A}}
\newcommand{\Mcal}{\mathcal{M}}
\newcommand{\Hcal}{\mathcal{H}}
\newcommand{\Dcal}{\mathcal{D}}
\newcommand{\Ncal}{\mathcal{N}}
\newcommand{\Ccal}{\mathcal{C}}
\newcommand{\Vcal}{\mathcal{V}}
\newcommand{\philam}{\phi_{\lambda\mu}}
\newcommand{\Vmat}{\mathbf{V}}
\newcommand{\Jset}{\mathcal{J}}
\newcommand{\dd}{\,\mathrm{d}}
\newcommand{\op}[1]{\operatorname{#1}}
\newcommand{\spec}{\sigma}
\newcommand{\Norm}[1]{\left\lVert #1 \right\rVert}
\newcommand{\Abs}[1]{\left\lvert #1 \right\rvert}
\newcommand{\ip}[2]{\left\langle #1, #2 \right\rangle}
\newcommand{\R}{\mathbb{R}}
\newcommand{\C}{\mathbb{C}}
\newcommand{\N}{\mathbb{N}}
\newcommand{\Z}{\mathbb{Z}}
\newcommand{\Q}{\mathbb{Q}}
\newcommand{\Pmark}{\mathrm{P}}
\newcommand{\Nmark}{\mathrm{N}}
\newcommand{\Omark}{\mathrm{O}}
\newcommand{\Fmark}{\mathrm{F}}

% ============================================================
% TITLE
% ============================================================
\title{%
  \textbf{A Multirepresentational Construction--Reconstruction\\
  Methodology for Mathematical Structures}\\[4pt]
  \large Part V: Validation Matrix, Worked Example, and Formalisation
}
\author{}
\date{}

\begin{document}
\maketitle

\begin{abstract}
\noindent
Part~V closes the methodology. We first consolidate the validation
matrix $\Vmat$ introduced in Part~I and populated in Parts~II--IV,
giving it a canonical form and a set of closure criteria. We then
present a fully worked example traversing all six construction
layers, the reconstruction pipeline, and the discrete-to-continuum
verification chain. Finally, we treat the formalisation of the
methodology in proof assistants, distinguishing the mathematics from
its mechanical encoding, and we state the limits of what can be
verified computationally. Appendices collect the notation table, the
main proofs, the numerical algorithms, and the bibliography.
\end{abstract}

\tableofcontents

% ============================================================
\part{Validation Matrix, Worked Example, and Formalisation}
% ============================================================

\section{The Validation Matrix}
\label{sec:validation-matrix}

\subsection{Recapitulation of the Validation Criteria}

Recall from Part~I, Definition~1.6, the set of validation criteria
\[
  \Jset = \{\mathrm{Ex}, \mathrm{Co}, \mathrm{In}, \mathrm{Nu},
          \mathrm{St}, \mathrm{Cv}, \mathrm{Re}, \mathrm{Rev}\},
\]
with the following interpretation:
\begin{center}
\begin{tabular}{c l}
  \toprule
  Symbol & Meaning \\
  \midrule
  $\mathrm{Ex}$  & Existence \\
  $\mathrm{Co}$  & Coherence (functoriality, naturality) \\
  $\mathrm{In}$  & Invariance under admissible transformations \\
  $\mathrm{Nu}$  & Numerical consistency \\
  $\mathrm{St}$  & Stability under perturbations \\
  $\mathrm{Cv}$  & Convergence as $N \to \infty$ \\
  $\mathrm{Re}$  & Regularity of the limit \\
  $\mathrm{Rev}$ & Reversibility of the construction \\
  \bottomrule
\end{tabular}
\end{center}

\begin{definition}[Epistemic statuses]\label{def:statuses}
The entries of $\Vmat$ take values in
\[
  \{\Pmark, \Nmark, \Omark, \Fmark\}
\]
with the following interpretation:
\begin{itemize}
  \item $\Pmark$: the criterion is \emph{proved} mathematically;
  \item $\Nmark$: the criterion is \emph{numerically supported} but
        not proved;
  \item $\Omark$: the criterion is \emph{open} at the time of
        writing;
  \item $\Fmark$: the criterion \emph{fails} for the given problem.
\end{itemize}
\end{definition}

\begin{remark}
The distinction between $\Pmark$ and $\Nmark$ is the operational
form of the distinction between mathematics and computation
developed throughout Parts~I--IV. A numerical result, however
stable, is not a proof; it is at best evidence for the existence of
a proof.
\end{remark}

\subsection{The Canonical Form of $\Vmat$}

\begin{definition}[Canonical validation matrix]\label{def:canonical-matrix}
The \emph{canonical validation matrix} of the methodology is
\[
  \Vmat
  \;=\;
  \bigl(v_{\lambda j}\bigr)_{\substack{\lambda \in \Lambda^\ast \\ j \in \Jset}},
\]
where
\[
  \Lambda^\ast
  \;=\;
  \{\Sp, \Ge, \To, \Al, \Nu, \Fo, \Cl, \Corr, \An\}
\]
is the extended layer set comprising the six construction layers and
the three reconstruction layers.
\end{definition}

\begin{protocol}[Population of $\Vmat$]\label{prot:population}
The entry $v_{\lambda j}$ is populated as follows:
\begin{enumerate}[label=(\roman*)]
  \item if the criterion $j$ has been established by a proof at layer
        $\lambda$, set $v_{\lambda j} = \Pmark$;
  \item if the criterion $j$ has been verified numerically at layer
        $\lambda$ but no proof is available, set
        $v_{\lambda j} = \Nmark$;
  \item if the criterion $j$ is neither proved nor numerically
        supported, and no counterexample is known, set
        $v_{\lambda j} = \Omark$;
  \item if the criterion $j$ is known to fail at layer $\lambda$,
        set $v_{\lambda j} = \Fmark$.
\end{enumerate}
\end{protocol}

\subsection{Instantiation}

The following table displays the instantiation of $\Vmat$ obtained
by combining the layerwise tables of Parts~II--IV.

\begin{center}
\renewcommand{\arraystretch}{1.15}
\begin{tabular}{l|cccccccc}
  \toprule
  Layer $\backslash$ Criterion
  & $\mathrm{Ex}$ & $\mathrm{Co}$ & $\mathrm{In}$ & $\mathrm{Nu}$
  & $\mathrm{St}$ & $\mathrm{Cv}$ & $\mathrm{Re}$ & $\mathrm{Rev}$ \\
  \midrule
  $\Sp$   & $\Pmark$ & $\Pmark$ & $\Pmark$ & $\Nmark$ & $\Pmark$ & $\Pmark$ & $\Omark$ & $\Omark$ \\
  $\Ge$   & $\Pmark$ & $\Pmark$ & $\Pmark$ & $\Nmark$ & $\Omark$ & $\Omark$ & $\Pmark$ & $\Omark$ \\
  $\To$   & $\Pmark$ & $\Pmark$ & $\Pmark$ & $\Nmark$ & $\Pmark$ & $\Pmark$ & $\Pmark$ & $\Omark$ \\
  $\Al$   & $\Pmark$ & $\Pmark$ & $\Pmark$ & $\Nmark$ & $\Omark$ & $\Omark$ & $\Pmark$ & $\Omark$ \\
  $\Nu$   & $\Pmark$ & $\Pmark$ & $\Pmark$ & $\Nmark$ & $\Pmark$ & $\Pmark$ & $\Omark$ & $\Omark$ \\
  $\Fo$   & $\Pmark$ & $\Pmark$ & $\Pmark$ & $\Omark$ & $\Pmark$ & $\Pmark$ & $\Pmark$ & $\Pmark$ \\
  $\Cl$   & $\Pmark$ & $\Pmark$ & $\Pmark$ & $\Omark$ & $\Omark$ & $\Omark$ & $\Pmark$ & $\Omark$ \\
  $\Corr$ & $\Omark$ & $\Omark$ & $\Omark$ & $\Omark$ & $\Omark$ & $\Pmark$ & $\Omark$ & $\Omark$ \\
  $\An$   & $\Omark$ & $\Omark$ & $\Omark$ & $\Omark$ & $\Omark$ & $\Omark$ & $\Pmark$ & $\Omark$ \\
  \bottomrule
\end{tabular}
\end{center}

\begin{remark}
The entries of this table are \emph{problem-dependent}: they reflect
the state of knowledge for a generic problem in which each layer is
handled by the standard methods cited in Parts~II--IV. Specific
problems will have specific instantiations.
\end{remark}

\subsection{Closure Criteria}

\begin{definition}[Row closure]\label{def:row-closure}
A layer $\lambda$ is \emph{closed} if
\[
  v_{\lambda j} \in \{\Pmark\} \quad \text{for at least one } j \in \Jset,
\]
and
\[
  v_{\lambda j} \neq \Fmark \quad \text{for all } j \in \Jset .
\]
\end{definition}

\begin{definition}[Full closure]\label{def:full-closure}
The methodology is \emph{fully closed} if every layer
$\lambda \in \Lambda^\ast$ is closed in the sense of
Definition~\ref{def:row-closure}, and the reversibility criterion
$\mathrm{Rev}$ is marked $\Pmark$ at the formal layer.
\end{definition}

\begin{theorem}[Closure of the methodology]\label{thm:closure-V}
If every layer is closed and the reversible criterion is marked
$\Pmark$ at the formal layer, then the methodology yields a
bidirectional characterisation
\[
  \Ob \;\rightleftarrows\; \Ecal_\Corr
\]
in the sense of Part~III, Theorem~3.12.
\end{theorem}

\begin{proof}
Combine the layerwise closure conditions with the closure criterion
of Part~III, Theorem~3.14, and the coherence of the construction
functor established in Part~I, Theorem~1.20.
\end{proof}

\subsection{Obstruction Analysis}

\begin{definition}[Obstruction]\label{def:obstruction}
An \emph{obstruction} is an entry $v_{\lambda j} = \Fmark$ in the
validation matrix. The \emph{obstruction layer} is the corresponding
$\lambda$, and the \emph{obstruction criterion} is the corresponding
$j$.
\end{definition}

\begin{proposition}[Localisation of obstructions]\label{prop:localisation}
Every obstruction is localised at a unique pair
$(\lambda, j) \in \Lambda^\ast \times \Jset$. Consequently,
obstructions can be enumerated, ranked, and resolved layer by layer.
\end{proposition}

\begin{proof}
Immediate from Definition~\ref{def:obstruction}.
\end{proof}

\begin{example}[Spectral obstruction]
The inverse spectral problem (Part~III, Example~3.6) exhibits an
obstruction at $(\Sp, \mathrm{Rev})$: injectivity of $\rho_\Inv$
fails because isospectral non-isometric domains exist
(Gordon--Webb--Wolpert). The obstruction is localised at the spectral
layer, criterion of reversibility.
\end{example}

\begin{example}[Numerical obstruction]
A discretisation scheme whose condition number grows unboundedly
with $N$ exhibits an obstruction at $(\Nu, \mathrm{St})$: stability
fails. The obstruction is localised at the numerical layer,
criterion of stability.
\end{example}

% ------------------------------------------------------------
\section{Worked Example}
\label{sec:worked-example}
% ------------------------------------------------------------

The following example illustrates the full methodology on a
concrete problem. The objective is not to solve a new research
problem but to display the traversal of every layer with explicit
content.

\subsection{Statement}

\begin{example}[Worked problem: the vibrating string]\label{ex:vibrating-string}
Let $\Omega = (0, \pi) \subset \R$, and consider the eigenvalue
problem
\[
  -u''(x) = \lambda u(x), \qquad x \in \Omega,
\]
with Dirichlet boundary conditions
\[
  u(0) = u(\pi) = 0 .
\]
The classical solution is
\[
  \lambda_n = n^2, \qquad u_n(x) = \sin(n x),
  \qquad n \in \N, \; n \geq 1 .
\]
\end{example}

\subsection{Layer $\Sp$: Spectral Datum}

\begin{construction}[Spectral construction]\label{constr:spectral-worked}
The spectral datum is
\[
  \Ocal_\Sp = (T, \spec(T), E_T)
\]
with $T = -\frac{d^2}{dx^2}$ on the domain
$D(T) = H^2(\Omega) \cap H^1_0(\Omega)$. The operator $T$ is
self-adjoint, positive, and unbounded. Its spectrum is
\[
  \spec(T) = \{n^2 : n \in \N, n \geq 1\},
\]
with multiplicities all equal to $1$. The spectral measure is
\[
  E_T(B) \varphi = \sum_{n \in \N, \, n^2 \in B} \ip{\varphi}{u_n} u_n,
  \qquad B \in \mathcal{B}(\spec(T)).
\]
\end{construction}

\begin{proposition}[Spectral validation]\label{prop:spectral-worked}
For the vibrating string:
\[
  v_{\Sp, \mathrm{Ex}} = \Pmark, \quad
  v_{\Sp, \mathrm{Co}} = \Pmark, \quad
  v_{\Sp, \mathrm{In}} = \Pmark, \quad
  v_{\Sp, \mathrm{Cv}} = \Pmark .
\]
\end{proposition}

\begin{proof}
Existence follows from the spectral theorem for unbounded
self-adjoint operators (Reed--Simon, Vol.~I, Theorem~VIII.6).
Coherence follows from functoriality of the spectral measure.
Invariance follows from unitary invariance of the spectrum
(Part~II, Proposition~2.6). Convergence follows from the Weyl
asymptotics
\[
  \lambda_n = n^2 + \mathcal{O}(1) .
\]
\end{proof}

\subsection{Layer $\Ge$: Geometric Datum}

\begin{construction}[Geometric construction]\label{constr:geom-worked}
The geometric realisation is
\[
  M = \mathbb{S}^1_{\pi} := \R / 2\pi \Z,
\]
the circle of circumference $2\pi$, with the standard Riemannian
metric
\[
  g = \dd x \otimes \dd x .
\]
The Levi-Civita connection is the flat connection
$\nabla = \partial_x$, and the Riemann curvature tensor vanishes:
\[
  R = 0 .
\]
The Laplace--Beltrami operator on $(M, g)$ is
\[
  \Delta_g = \frac{d^2}{dx^2} = -T .
\]
\end{construction}

\subsection{Layer $\To$: Topological Datum}

\begin{construction}[Topological construction]\label{constr:top-worked}
The underlying topological space is $X = M = \mathbb{S}^1$. Its
homology and homotopy are
\[
  H_0(X; \Z) = \Z, \quad H_1(X; \Z) = \Z, \quad H_k(X; \Z) = 0
  \quad (k \geq 2),
\]
\[
  \pi_1(X) = \Z, \quad \pi_k(X) = 0 \quad (k \geq 2),
\]
\[
  \chi(X) = 0 .
\]
\end{construction}

\subsection{Layer $\Al$: Algebraic Datum}

\begin{construction}[Algebraic construction]\label{constr:alg-worked}
The algebraic datum is
\[
  \Ocal_\Al = (A, G, R, \Mcal)
\]
with:
\begin{itemize}
  \item $A = \C[\partial_x]$, the polynomial algebra in the
        derivation $\partial_x$;
  \item $G = \mathbb{Z}$, the fundamental group of $X$ acting by
        rotations;
  \item $R = \Z$, the ring of integers encoding the arithmetic of
        the spectrum;
  \item $\Mcal = C^\infty(X) \otimes \C$, the module of smooth
        functions.
\end{itemize}
The eigenvalue equation is
\[
  (\partial_x^2 + \lambda) u = 0,
\]
which identifies $A$ as a quotient of $\C[\partial_x]$.
\end{construction}

\subsection{Layer $\Nu$: Numerical Datum}

\begin{construction}[Numerical construction]\label{constr:num-worked}
Choose the finite-difference discretisation with step
$h = \pi / (N+1)$ and grid points
\[
  x_i = i h, \qquad i = 0, 1, \ldots, N+1 .
\]
The discrete operator is the tridiagonal matrix
\[
  T_N = \frac{1}{h^2}
  \begin{pmatrix}
    2 & -1 & & \\
    -1 & 2 & -1 & \\
    & \ddots & \ddots & \ddots \\
    & & -1 & 2
  \end{pmatrix}
  \in \mathrm{Mat}_{N \times N}(\R),
\]
acting on $u = (u_1, \ldots, u_N)^\top$ with
$u_0 = u_{N+1} = 0$.
\end{construction}

\begin{proposition}[Eigenvalue approximation]\label{prop:eigen-approx}
The eigenvalues of $T_N$ are
\[
  \lambda_n^{(N)}
  \;=\;
  \frac{4}{h^2} \sin^2\!\left( \frac{n \pi h}{2} \right),
  \qquad n = 1, \ldots, N,
\]
and satisfy
\[
  \Abs{\lambda_n^{(N)} - n^2}
  \;\leq\;
  \frac{n^4 h^2}{12} + \mathcal{O}(h^4) .
\]
In particular,
\[
  \lambda_n^{(N)} \to n^2 \qquad (N \to \infty).
\]
\end{proposition}

\begin{proof}
Standard computation: the eigenvalue equation reduces to a linear
recurrence with characteristic roots $e^{\pm i \theta}$, where
$\theta$ satisfies $\cos \theta = 1 - h^2 \lambda / 2$. Solving for
$\lambda$ gives the displayed expression; the estimate follows from
Taylor expansion of $\sin^2$ around the origin.
\end{proof}

\subsection{Layer $\Fo$: Formal Datum}

\begin{construction}[Formal construction]\label{constr:formal-worked}
The formal language is $\mathcal{L} = \mathcal{L}_{\mathrm{PDE}}$,
a first-order language with sorts for $\R$, functions
$\partial_x, \partial_x^2, u$, and a constant symbol $\lambda$. The
axioms are
\[
  \Gamma = \{\;
    \forall x \in \Omega,\; u(0) = u(\pi) = 0, \;
    \forall x \in \Omega,\; -u''(x) = \lambda u(x) \;
  \}.
\]
Derivability is given by natural deduction. The completeness theorem
(Part~II, Theorem~7.9) provides
\[
  \Gamma \vdash \varphi
  \;\Longleftrightarrow\;
  \Gamma \models \varphi .
\]
\end{construction}

\subsection{Reconstruction: Classical Extraction}

\begin{construction}[Reconstruction]\label{constr:reconstruct-worked}
The reconstruction $\rho_\Cl$ extracts the classical model
\[
  \Mcal_\star = \bigl( \Omega, \; -u'' = \lambda u \bigr)
\]
as the distinguished classical equation. A direct comparison with
the classical literature identifies
\[
  \mathcal{E}_\Cl = \{-u'' = \lambda u\}
  \;\trianglelefteq\;
  \Ecal_\Cl ,
\]
where $\trianglelefteq$ is the containment relation of Part~III,
Definition~3.3. Consequently, by Part~III,
Proposition~3.4,
\[
  \Delta L = 0 .
\]
\end{construction}

\begin{proposition}[No correction]\label{prop:no-correction}
For the vibrating string, the delta operator vanishes identically:
\[
  \Delta L = 0, \qquad L_\Corr = L_\Cl .
\]
\end{proposition}

\begin{proof}
By Construction~\ref{constr:reconstruct-worked} and Part~III,
Proposition~3.4.
\end{proof}

\subsection{Discrete-to-Continuum Passage}

\begin{proposition}[Convergence of the discrete family]\label{prop:convergence-worked}
The discrete family $\{T_N\}_N$ satisfies the three-proof structure
of Part~IV, Definition~4.12:
\[
  v_{\Nu, \mathrm{St}} = \Pmark, \quad
  v_{\Nu, \mathrm{Cv}} = \Pmark, \quad
  v_{\Nu, \mathrm{Nu}} = \Nmark .
\]
\end{proposition}

\begin{proof}
Consistency: $T_N \pi_N u \to \pi_N T u$ by Taylor expansion.
Stability: $T_N$ is diagonally dominant and positive definite, so
$\Norm{T_N^{-1}} \leq C$ uniformly in $N$ (Gershgorin). Convergence
follows from Theorem~4.6 (Lax--Richtmyer analogue).
\end{proof}

\subsection{Analytic Closure}

\begin{proposition}[Analyticity]\label{prop:analytic-worked}
Every eigenvalue $\lambda_n = n^2$ admits a corresponding
eigenfunction $u_n(x) = \sin(n x)$ that is real-analytic on
$\Omega = (0, \pi)$.
\end{proposition}

\begin{proof}
The function $\sin(n x)$ extends to an entire function on $\C$ and
satisfies the ODE $-u'' = n^2 u$ identically. Analyticity follows
from the analytic regularity theorem (Part~IV,
Theorem~4.15).
\end{proof}

\subsection{Summary of the Worked Example}

The following table summarises the validation status for the
vibrating string problem.

\begin{center}
\begin{tabular}{l|cccccccc}
  \toprule
  Layer & $\mathrm{Ex}$ & $\mathrm{Co}$ & $\mathrm{In}$ & $\mathrm{Nu}$
  & $\mathrm{St}$ & $\mathrm{Cv}$ & $\mathrm{Re}$ & $\mathrm{Rev}$ \\
  \midrule
  $\Sp$   & $\Pmark$ & $\Pmark$ & $\Pmark$ & $\Nmark$ & $\Pmark$ & $\Pmark$ & $\Pmark$ & $\Pmark$ \\
  $\Ge$   & $\Pmark$ & $\Pmark$ & $\Pmark$ & $\Nmark$ & $\Pmark$ & $\Pmark$ & $\Pmark$ & $\Pmark$ \\
  $\To$   & $\Pmark$ & $\Pmark$ & $\Pmark$ & $\Omark$ & $\Pmark$ & $\Pmark$ & $\Pmark$ & $\Pmark$ \\
  $\Al$   & $\Pmark$ & $\Pmark$ & $\Pmark$ & $\Omark$ & $\Pmark$ & $\Pmark$ & $\Pmark$ & $\Pmark$ \\
  $\Nu$   & $\Pmark$ & $\Pmark$ & $\Pmark$ & $\Nmark$ & $\Pmark$ & $\Pmark$ & $\Omark$ & $\Omark$ \\
  $\Fo$   & $\Pmark$ & $\Pmark$ & $\Pmark$ & $\Omark$ & $\Pmark$ & $\Pmark$ & $\Pmark$ & $\Pmark$ \\
  $\Cl$   & $\Pmark$ & $\Pmark$ & $\Pmark$ & $\Omark$ & $\Pmark$ & $\Pmark$ & $\Pmark$ & $\Pmark$ \\
  $\Corr$ & $\Pmark$ & $\Pmark$ & $\Pmark$ & $\Omark$ & $\Pmark$ & $\Pmark$ & $\Pmark$ & $\Pmark$ \\
  $\An$   & $\Pmark$ & $\Pmark$ & $\Pmark$ & $\Omark$ & $\Pmark$ & $\Pmark$ & $\Pmark$ & $\Pmark$ \\
  \bottomrule
\end{tabular}
\end{center}

\begin{remark}
The vibrating string problem admits a \emph{fully closed}
methodology: every layer satisfies every applicable criterion either
by proof or, in the case of numerical criteria, by numerically
supported evidence. This is the ideal case; it is not the generic
case.
\end{remark}

% ------------------------------------------------------------
\section{Formalisation in Proof Assistants}
\label{sec:proof-assistants}
% ------------------------------------------------------------

\subsection{The Two Notions of Formalisation}

\begin{definition}[Mathematical formalisation]\label{def:math-formalisation}
\emph{Mathematical formalisation} is the process of expressing a
mathematical argument as a sequence of statements and inferences in
a fixed language $\mathcal{L}$, using a fixed set of axioms
$\Gamma$, and a fixed set of inference rules.
\end{definition}

\begin{definition}[Mechanical formalisation]\label{def:mech-formalisation}
\emph{Mechanical formalisation} is the process of encoding a
mathematical formalisation in a proof assistant (Lean, Coq, Isabelle,
Agda, \ldots) so that the derivability relation $\vdash$ is checked
by a program.
\end{definition}

\begin{remark}
Mathematical formalisation does not require mechanical formalisation.
Every rigorous proof in the mathematical literature is a
mathematical formalisation; only a small fraction has been encoded
mechanically. The distinction parallels, at the level of proofs,
the distinction between mathematical truth and numerical evidence
developed throughout Parts~I--IV.
\end{remark}

\subsection{Encoding the Metamodel}

\begin{construction}[Encoding the master object]\label{constr:encoding-Ob}
In Lean~4 with Mathlib, the master object $\Ob$ and its layers can
be encoded as a dependent record:
\begin{verbatim}
structure MasterObject where
  Sp : SpectralDatum
  Ge : GeometricDatum
  To : TopologicalDatum
  Al : AlgebraicDatum
  Nu : NumericalDatum
  Fo : FormalDatum
\end{verbatim}
Each layer is a structure with its primitive fields and its
invariants. The transition morphisms are functions between these
structures, and the functoriality axiom is a proof obligation.
\end{construction}

\begin{construction}[Encoding the transitions]\label{constr:encoding-transitions}
The transition morphisms of Part~II, \S2.4--\S7.4 become functions:
\begin{verbatim}
def phi_Sp_Ge : SpectralDatum → GeometricDatum := ...
def phi_Ge_To : GeometricDatum → TopologicalDatum := ...
def phi_To_Al : TopologicalDatum → AlgebraicDatum := ...
def phi_Al_Nu : AlgebraicDatum → NumericalDatum := ...
def phi_Nu_Fo : NumericalDatum → FormalDatum := ...
\end{verbatim}
Functoriality (Part~I, Axiom~1.3) becomes a lemma:
\begin{verbatim}
lemma phi_functorial
  (h : phi_Ge_To ∘ phi_Sp_Ge = phi_Sp_To) : ... := ...
\end{verbatim}
\end{construction}

\subsection{Verification Obligations}

\begin{protocol}[Mechanical verification obligations]\label{prot:mech-obligations}
For the methodology to be mechanically verified, the following
obligations must be discharged in the proof assistant:
\begin{enumerate}[label=(\roman*)]
  \item \textbf{Functoriality:} each layer transition respects
        composition (Part~I, Axiom~1.3).
  \item \textbf{Projection coherence:} the discrete projections
        commute with the continuum projections (Part~I,
        Axiom~1.4).
  \item \textbf{Correspondence limit:} the corrected operator admits
        the classical operator as a limit (Part~III,
        Axiom~3.4).
  \item \textbf{Closure:} the bidirectional characterisation
        $\Ob \rightleftarrows \Ecal_\Corr$ is proved
        (Part~III, Theorem~3.12).
  \item \textbf{Chain closure:} the final verification chain
        terminates in an analytic solution (Part~IV,
        Theorem~4.23).
\end{enumerate}
\end{protocol}

\begin{theorem}[Mechanical verifiability]\label{thm:mech-verifiability}
If the obligations (i)--(v) of Protocol~\ref{prot:mech-obligations}
are discharged in a consistent proof assistant, then the
methodology is mechanically verified.
\end{theorem}

\begin{proof}
Each obligation is a formal statement in the language of the proof
assistant. By the soundness of the assistant's logic (a standard
theorem for Lean, Coq, Isabelle, and Agda), each discharged
obligation is a theorem in the underlying formal system. Their
conjunction is the mechanical verification of the methodology.
\end{proof}

\subsection{Limits of Mechanical Verification}

\begin{proposition}[Non-verifiable components]\label{prop:non-verif}
The following components of the methodology \emph{cannot} be
mechanically verified in isolation:
\begin{enumerate}[label=(\alph*)]
  \item the choice of the distinguished classical model
        $\Mcal_\star$ (Part~III, Construction~3.2);
  \item the choice of the discretisation scheme
        (Part~IV, Definition~4.1);
  \item the choice of the detection algorithm
        (Part~IV, Definition~4.11);
  \item the interpretation of the formal language in the intended
        mathematical universe (Part~II, Construction~7.6).
\end{enumerate}
\end{proposition}

\begin{proof}
Each of (a)--(d) is a \emph{choice} rather than a \emph{theorem}: it
is an input to the methodology, not a consequence of it. Mechanical
verification can check that the consequences of a choice are
correct, but not that the choice itself is correct.
\end{proof}

\begin{remark}
The limits of mechanical verification mirror the limits of
mathematical formalism itself: a formal system can verify that
certain statements follow from certain axioms, but it cannot, from
within, verify that the axioms capture the intended mathematical
reality. This is the formal analogue of the distinction between
syntax and semantics, and it is why the validation matrix $\Vmat$
records both the proof status ($\Pmark$) and the numerical status
($\Nmark$) of each criterion.
\end{remark}

\subsection{Recommended Encodings}

\begin{example}[Lean~4 encoding of the spectral layer]
The spectral layer of Construction~\ref{constr:encoding-Ob} can be
refined as follows.
\begin{verbatim}
structure SpectralDatum (H : Type*) [NormedAddCommGroup H]
  [InnerProductSpace ℂ H] [CompleteSpace H] where
  T : H →ₗ.[ℂ] H
  closed : IsClosed (T.graph)
  selfAdjoint : IsSelfAdjoint T
  spectrum : Set ℂ
  spectrum_eq : spectrum = spectrum ℂ T
  spectralMeasure : ...
\end{verbatim}
The spectral theorem is available in Mathlib as
\texttt{IsSelfAdjoint.spectralTheorem}, and the correspondence
$\spec(T) \subset \R$ is the lemma
\texttt{IsSelfAdjoint.spectrum\_subset\_real}.
\end{example}

\begin{example}[Lean~4 encoding of the topological layer]
The topological layer can be encoded using Mathlib's
\texttt{SingularHomology} and \texttt{FundamentalGroup}.
\begin{verbatim}
structure TopologicalDatum where
  X : Type*
  topX : TopologicalSpace X
  H : ℕ → Type*
  H_eq : ∀ k, H k = SingularHomology X k
  pi1 : FundamentalGroup X
  chi : ℤ
  chi_eq : chi = eulerCharacteristic X
\end{verbatim}
Homotopy invariance is the lemma
\texttt{SingularHomology.homotopyInvariant}.
\end{example}

\begin{remark}
The two encodings above illustrate a general pattern: each layer of
the methodology can be encoded using existing Mathlib structures,
and each transition morphism becomes a proof obligation rather than
an object of study. This is the sense in which mechanical
formalisation is a \emph{consumer} of the methodology, not a
\emph{producer} of it.
\end{remark}

% ------------------------------------------------------------
\section*{Summary of Part V}
% ------------------------------------------------------------

Part~V has closed the methodology by:
\begin{itemize}
  \item consolidating the validation matrix $\Vmat$ in its
        canonical form (Definition~\ref{def:canonical-matrix});
  \item defining the closure criteria for individual layers
        (Definition~\ref{def:row-closure}) and for the methodology
        as a whole (Definition~\ref{def:full-closure});
  \item proving the closure theorem
        (Theorem~\ref{thm:closure-V});
  \item presenting a fully worked example (the vibrating string)
        traversing all nine layers;
  \item distinguishing mathematical formalisation from mechanical
        formalisation (Definitions~\ref{def:math-formalisation},
        \ref{def:mech-formalisation});
  \item specifying the mechanical verification obligations
        (Protocol~\ref{prot:mech-obligations});
  \item identifying the components that cannot be mechanically
        verified (Proposition~\ref{prop:non-verif}).
\end{itemize}

The central structural results are:
\[
  \text{every layer closed} \;\Longrightarrow\;
  \Ob \rightleftarrows \Ecal_\Corr ,
\]
\[
  \text{obligations (i)--(v) discharged} \;\Longrightarrow\;
  \text{mechanical verification},
\]
and the fundamental distinction
\[
  \Pmark \;\neq\; \Nmark .
\]

% ============================================================
\appendix
\part*{Appendices}
\addcontentsline{toc}{part}{Appendices}
% ============================================================

\section{Notation Table}
\label{app:notation}

\begin{center}
\renewcommand{\arraystretch}{1.2}
\begin{longtable}{c l}
  \toprule
  Symbol & Meaning \\
  \midrule
  $\Ob$                        & Master object \\
  $\Lam$                       & Layer index set \\
  $\Ocal_\lambda$              & Layer-$\lambda$ datum \\
  $\philam$                    & Transition morphism \\
  $\rho_\Cl, \rho_\Inv, \rho_\An$ & Reconstruction operators \\
  $\Ecal_\Cl$                  & Classical-model layer \\
  $\Ecal_\Corr$                & Corrected-model layer \\
  $\Acal$                      & Analytic layer \\
  $\Ob^{(N)}$                  & Discrete master object at resolution $N$ \\
  $\varepsilon_N$              & Layer-aggregated error \\
  $\Vmat$                      & Validation matrix \\
  $\Jset$                      & Set of validation criteria \\
  $\Pmark, \Nmark, \Omark, \Fmark$ & Proved, numerical, open, failed \\
  $\Delta L$                   & Delta operator (correction) \\
  $L_\Corr$                    & Corrected operator \\
  $\spec(T)$                   & Spectrum of $T$ \\
  $E_T$                        & Spectral measure of $T$ \\
  $\Delta_g$                   & Laplace--Beltrami operator \\
  $H_k(X)$                     & $k$-th singular homology \\
  $\pi_k(X)$                   & $k$-th homotopy group \\
  $\chi(X)$                    & Euler characteristic \\
  $\kappa(T_N)$                & Condition number \\
  \bottomrule
\end{longtable}
\end{center}

\section{Proofs}
\label{app:proofs}

\begin{proof}[Proof of Proposition~2.6 (unitary invariance)]
Let $U : H \to H'$ be unitary and $T' = U T U^{-1}$. If
$T v = \lambda v$, then $T' (U v) = U T v = \lambda U v$, so
$\spec(T') \supseteq \spec(T)$; the reverse inclusion follows by
symmetry. For the spectral measure: for $B \in \mathcal{B}(\spec(T))$,
\[
  E_{T'}(B)
  = \int_B \dd E_{T'}
  = U \left( \int_B \dd E_T \right) U^{-1}
  = U E_T(B) U^{-1}.
\]
\end{proof}

\begin{proof}[Proof of Theorem~4.6 (consistency--stability--convergence)]
The discrete equation is $\mathcal{F}_\lambda^{(N)} u_N = 0$.
Subtracting $\mathcal{F}_\lambda^{(N)} \pi_\lambda^{(N)} u$:
\[
  \mathcal{F}_\lambda^{(N)}(u_N - \pi_\lambda^{(N)} u)
  = -\mathcal{F}_\lambda^{(N)} \pi_\lambda^{(N)} u .
\]
Applying the inverse estimate:
\[
  \Norm{u_N - \pi_\lambda^{(N)} u}
  \leq \Norm{(\mathcal{F}_\lambda^{(N)})^{-1}}
        \Norm{\mathcal{F}_\lambda^{(N)} \pi_\lambda^{(N)} u}
  \leq C \Norm{\mathcal{F}_\lambda^{(N)} \pi_\lambda^{(N)} u} .
\]
The last factor is controlled by consistency
(Definition~4.5), giving the stated estimate.
\end{proof}

\begin{proof}[Proof of Theorem~4.19 (compactness)]
By Assumption~4.14, $\{u_N\}$ is bounded in $X$. Banach--Alaoglu
gives a weakly convergent subsequence $u_{N_j} \rightharpoonup u_\infty$
in $X$. The compact embedding $X \hookrightarrow Y$ upgrades weak
convergence in $X$ to strong convergence in $Y$.
\end{proof}

\section{Numerical Algorithms}
\label{app:algorithms}

\begin{protocol}[Eigenvalue computation]
Given $T_N \in \mathrm{Mat}_{N \times N}(\R)$ symmetric positive
definite:
\begin{enumerate}[label=(\arabic*)]
  \item Compute the Cholesky factorisation $T_N = L L^\top$.
  \item Apply the Lanczos algorithm with $k$ steps to obtain a
        tridiagonal matrix $T_k^{(L)}$ with
        $T_k^{(L)} \approx Q_k^\top T_N Q_k$.
  \item Compute the eigenvalues of $T_k^{(L)}$ by the QR algorithm.
  \item Return the $k$ Ritz values as approximations to the $k$
        smallest eigenvalues of $T_N$.
\end{enumerate}
\end{protocol}

\begin{protocol}[Persistent homology computation]
Given a filtered simplicial complex $X_1 \subset \cdots \subset X_N$:
\begin{enumerate}[label=(\arabic*)]
  \item Construct the boundary matrix $D$ of the full complex $X_N$.
  \item Reduce $D$ by column operations (standard algorithm) to
        obtain the reduced matrix $R$.
  \item Extract the barcode from the pivot positions of $R$.
  \item Persistence intervals are given by
        $(b_i, d_i) = (\text{birth}, \text{death})$.
\end{enumerate}
\end{protocol}

\begin{protocol}[Richardson extrapolation]
Given a sequence $\{u_N\}_N$ with asymptotic expansion
$u_N = u_\infty + a_1 N^{-p} + a_2 N^{-2p} + \cdots$:
\begin{enumerate}[label=(\arabic*)]
  \item Choose a base step $N_0$ and refinement factors
        $N_j = 2^j N_0$.
  \item Compute $u_{N_j}$ for $j = 0, 1, \ldots, q$.
  \item Solve the Vandermonde system for the coefficients
        $c_j^{(q)}$.
  \item Return $u^{(q)} = \sum_j c_j^{(q)} u_{N_j}$.
\end{enumerate}
\end{protocol}

\section{Bibliography}
\label{app:biblio}

\begin{thebibliography}{99}

\bibitem{reed-simon-1}
M.~Reed, B.~Simon,
\emph{Methods of Modern Mathematical Physics, Vol.~I: Functional
Analysis},
Academic Press, 1980.

\bibitem{reed-simon-4}
M.~Reed, B.~Simon,
\emph{Methods of Modern Mathematical Physics, Vol.~IV: Analysis of
Operators},
Academic Press, 1978.

\bibitem{ebbinghaus}
H.-D.~Ebbinghaus, J.~Flum, W.~Thomas,
\emph{Mathematical Logic},
Graduate Texts in Mathematics 291, Springer, 2021.

\bibitem{engl}
H.~W.~Engl, M.~Hanke, A.~Neubauer,
\emph{Regularization of Inverse Problems},
Kluwer, 1996.

\bibitem{morrey}
C.~B.~Morrey,
\emph{Multiple Integrals in the Calculus of Variations},
Springer, 1966.

\bibitem{bers-john}
L.~Bers, F.~John, M.~Schechter,
\emph{Partial Differential Equations},
AMS, 1964.

\bibitem{agmon}
S.~Agmon, A.~Douglis, L.~Nirenberg,
Estimates near the boundary for solutions of elliptic partial
differential equations,
\emph{Comm.~Pure Appl.~Math.}~12 (1959), 623--727.

\bibitem{gordon}
C.~Gordon, D.~Webb, S.~Wolpert,
Isospectral plane domains and surfaces via Riemannian orbifolds,
\emph{Invent.~Math.}~110 (1992), 1--22.

\bibitem{kac}
M.~Kac,
Can one hear the shape of a drum?
\emph{Amer.~Math.~Monthly}~73 (1966), 1--23.

\bibitem{lax-rich}
P.~D.~Lax, R.~D.~Richtmyer,
Survey of the stability of linear finite difference equations,
\emph{Comm.~Pure Appl.~Math.}~9 (1956), 267--293.

\bibitem{mathlib}
The Mathlib Community,
\emph{Mathlib: A unified library of mathematics formalized in Lean},
\url{https://leanprover-community.github.io/mathlib4_docs/}, 2024.

\end{thebibliography}

% ============================================================
\section*{Closing Remark}
% ============================================================

The methodology developed in Parts~I--V is best described as a
\emph{multirepresentational construction--reconstruction
methodology}. Its distinguishing features are:

\begin{enumerate}[label=(\roman*)]
  \item the treatment of mathematical representations as
        \emph{projections} of a single master object, rather than
        as distinct objects;
  \item the explicit separation of \emph{construction} (Part~II)
        from \emph{reconstruction} (Part~III);
  \item the treatment of the discrete-to-continuum passage
        (Part~IV) as a mathematical problem in its own right, with
        consistency, stability, convergence, and regularity as
        separate obligations;
  \item the systematic use of the validation matrix $\Vmat$ to
        record the epistemic status of each layer and criterion;
  \item the sharp distinction between
        \emph{mathematical} formalisation and \emph{mechanical}
        formalisation (Part~V, \S\ref{sec:proof-assistants}).
\end{enumerate}

The methodology does not produce mathematical results; it organises
them. Its value is diagnostic: it forces every representation of a
mathematical object to declare what it preserves, what it loses,
what it proves, and what it merely suggests. In this sense it is
closer to a \emph{protocol of mathematical hygiene} than to a
\emph{discovery method}.

\end{document}